\documentclass[11pt]{report}

\usepackage[margin=1in]{geometry}
\usepackage[T1]{fontenc}
\usepackage[utf8]{inputenc}
\usepackage{lmodern}
\usepackage{amsmath,amssymb,amsthm,bm,mathtools}
\usepackage{booktabs}
\usepackage{longtable}
\usepackage{array}
\usepackage{enumitem}
\usepackage{algorithm}
\usepackage{algpseudocode}
\usepackage{placeins}
\usepackage{xcolor}
\usepackage{hyperref}
\usepackage{url}

\hypersetup{
  colorlinks=true,
  linkcolor=blue!50!black,
  urlcolor=blue!50!black
}

\newtheorem{definition}{Definition}[chapter]
\newtheorem{remark}{Remark}[chapter]
\newtheorem{proposition}{Proposition}[chapter]

\newcommand{\R}{\mathbb{R}}
\newcommand{\E}{\mathbb{E}}
\newcommand{\Prob}{\mathbb{P}}
\newcommand{\calP}{\mathcal{P}}
\newcommand{\KL}{\operatorname{KL}}
\newcommand{\W}{\operatorname{W}}
\newcommand{\diag}{\operatorname{diag}}
\newcommand{\argmin}{\operatorname*{arg\,min}}
\newcommand{\argmax}{\operatorname*{arg\,max}}
\newcommand{\pushf}[2]{\left(#1\right)_{\#} #2}
\newcommand{\code}[1]{\texttt{#1}}

\title{Recent Neural Optimal Transport Candidates for BayesFilter\\
A Teaching Note for First-Year PhD Students}
\author{BayesFilter literature reread note}
\date{2026-06-22}

\begin{document}
\pagenumbering{roman}
\maketitle
\tableofcontents
\clearpage
\pagenumbering{arabic}

\begin{abstract}
This note is not written as a compressed survey.  It is written as a teaching
note for first-year PhD students who want to understand what recent
neural-optimal-transport papers are actually computing and why BayesFilter might
care.  The organizing principle is simple: a paper may learn a coupling, a warm
start, a static transport map, a conditional posterior map, a dynamic path, or a
measure-to-measure operator.  Those objects are not interchangeable.  Before one
can rank papers, one must understand what object each paper learns, how that
object is defined mathematically, and what would have to be deployed inside
BayesFilter.

The note therefore does two things at once.  First, it reproduces the central
mathematical objects and equations from the papers.  Second, it explicitly fills
in explanatory gaps when the original papers jump steps, move too quickly, or
leave important implications implicit.  The main teaching line focuses on a small
set of anchor papers: Makkuva et al. (2020) for the OT-ICNN baseline, Wang et
al. (2025) for conditional posterior transport, Al-Jarrah et al. (2025) for
filtering-native conditional transport, Hosseini et al. (2025) for theorem-level
conditional OT, and Meta OT / UNOT for retained-teacher acceleration.  The
broader shelf remains present, but it is treated later and more compactly.

One correction should be stated at the outset.  BayesFilter's earlier negative
``OT-ICNN'' result does not show that the original Makkuva et al. method fails.
The earlier benchmark tested bridge-style surrogates, not a faithful
reproduction of the paper's paired-convex-potential minimax method.  So the
right lesson is not ``OT-ICNN was disproven,'' but rather ``that benchmark did
not justify continued investment in that particular bridge route, and newer
conditional or retained-teacher papers are now more central.''
\end{abstract}

\chapter{Primer: what BayesFilter is trying to compute}

\section{A weighted cloud is not yet a transport rule}

At a filtering step, BayesFilter often has a weighted empirical measure
\begin{equation}
\mu_N^w = \sum_{i=1}^{N} w_i\delta_{x_i},
\qquad
w_i\ge 0,
\qquad
\sum_{i=1}^{N} w_i=1.
\label{eq:weighted-cloud}
\end{equation}
This is only a distributional summary.  It is not yet a transport rule.  To see
what is missing, start with a tiny three-particle example:
\begin{equation}
\mu_3^w = \tfrac12 \delta_{-1} + \tfrac13 \delta_{0} + \tfrac16 \delta_{2}.
\label{eq:three-particle-source}
\end{equation}
Suppose we want an equal-weight cloud
\begin{equation}
\bar\mu_3 = \tfrac13 \delta_{z_1} + \tfrac13 \delta_{z_2} + \tfrac13 \delta_{z_3}.
\label{eq:three-particle-target}
\end{equation}
A scalar OT cost does not tell us what the new points $z_j$ are.  What we need is
one of the following:
\begin{enumerate}[label=(\roman*)]
  \item a \emph{coupling} $\pi$ between source and target weights,
  \item a \emph{map} $T$ sending source points to target points,
  \item or a \emph{conditional update rule} that depends on an observation $y$.
\end{enumerate}

In the coupling view, one solves over matrices
\begin{equation}
\Pi(w,\bar w)
=
\left\{
\pi\in\R_+^{N\times N}:
\pi\mathbf 1 = w,
\ \pi^{\top}\mathbf 1 = \bar w
\right\},
\label{eq:transport-polytope}
\end{equation}
and then forms the equal-weight points by barycentric projection,
\begin{equation}
z_j = \frac{1}{\bar w_j}\sum_{i=1}^{N} \pi_{ij}x_i.
\label{eq:barycentric-projection}
\end{equation}
In the direct-map view, one instead learns a map $T$ and sets
\begin{equation}
z_i = T(x_i).
\label{eq:direct-map-deployment}
\end{equation}
In the conditional-map view, one learns a family of maps $T_y$ and sets
\begin{equation}
z_i = T_y(x_i).
\label{eq:conditional-map-deployment}
\end{equation}
This small example already teaches the central lesson of the whole note: two
papers can both say ``optimal transport'' while learning completely different
objects.

\paragraph{Implementation safety warning.}
A weighted cloud by itself is \emph{not} yet something one can deploy as an update
algorithm.  Before implementing any paper in the note, a student should force one
classification:
\begin{enumerate}[label=(\roman*)]
  \item is the paper learning a coupling or plan,
  \item a direct map,
  \item a conditional update family,
  \item a path or operator,
  \item or only a latent state that helps a retained teacher?
\end{enumerate}
Many wrong implementations come from skipping that classification and jumping
straight from a weighted empirical measure to code.

\section{A running conditional example}

The next step is to add conditioning.  Suppose
\begin{equation}
Z \sim \mathcal N(0,1),
\qquad
X\mid Y=y \sim \mathcal N(m(y),s(y)^2).
\label{eq:running-conditional-law}
\end{equation}
Then one perfectly valid conditional transport map is
\begin{equation}
g(z;y) = m(y) + s(y)z.
\label{eq:running-conditional-map}
\end{equation}
This one-line example is extremely important.  It tells us that a conditional
transport paper is trying to learn something like ``for each fixed observation
$y$, push a simple reference law to the posterior law.''  The inverse map is
\begin{equation}
g^{-1}(x;y)=\frac{x-m(y)}{s(y)}.
\label{eq:running-conditional-inverse}
\end{equation}
Therefore, by change of variables,
\begin{equation}
p_{X\mid Y}(x\mid y)
=
\varphi\!\bigl(g^{-1}(x;y)\bigr)
\left|\partial_x g^{-1}(x;y)\right|,
\label{eq:running-change-of-variables}
\end{equation}
where $\varphi$ is the standard Gaussian density.  Many of the papers below are
more complicated versions of exactly this story.

\section{Three object classes}

\paragraph{Retained-teacher methods.}
A retained-teacher method keeps a declared OT or entropic-OT solve and learns only
a warm start or latent solver state,
\begin{equation}
h_\theta(\mu,\nu)
\leadsto
\widehat \pi_\theta
\xrightarrow{\text{corrective solve}}
\pi_{\mathrm{teacher}}.
\label{eq:retained-teacher-schema}
\end{equation}
The learned object is not trusted as the final answer.  It only helps the teacher.

\paragraph{Direct static or conditional maps.}
A direct-map paper tries to learn
\begin{equation}
T_\theta : \R^d \to \R^d,
\qquad
\pushf{T_\theta}{\mu} \approx \nu,
\label{eq:direct-map-schema}
\end{equation}
or, in the conditional case,
\begin{equation}
\pushf{T_\theta(\cdot;y)}{\eta} \approx \pi(\cdot\mid y).
\label{eq:conditional-map-schema}
\end{equation}
Now the learned object itself is deployed.

\paragraph{Dynamic paths and measure operators.}
A richer family learns a path or operator.  A dynamic path satisfies a continuity
law,
\begin{equation}
\partial_t \rho_t + \nabla\cdot(\rho_t v_t)=0,
\qquad
t\in[0,1],
\label{eq:continuity-path}
\end{equation}
while an operator learns a mapping
\begin{equation}
\mathcal G_\theta : \calP(\R^d) \to \calP(\R^d).
\label{eq:measure-operator}
\end{equation}
These methods are often scientifically interesting, but they are usually farther
from BayesFilter's present deployment problem.

\paragraph{Implementation-safe object registry.}
For the rest of the note, keep four labels separate:
\begin{enumerate}[label=(\roman*)]
  \item \textbf{learned object} --- what the network or optimizer is actually fitting,
  \item \textbf{optimized object} --- what the loss or solver updates directly,
  \item \textbf{auxiliary object} --- critic, dual, value function, PDE residual, or solver state used only to enforce the contract,
  \item \textbf{deployed object} --- the thing BayesFilter would actually call at inference time.
\end{enumerate}
If a student cannot fill in those four labels for a paper, they should not start
coding yet.  That missing classification is itself the first bug signal.

\section{House notation and how to read the rest}

The note uses one house notation system throughout:
\begin{itemize}
  \item $\mu$ is always a source measure and $\nu$ a target measure unless a paper
  is intrinsically conditional, in which case the target is often written as
  $\pi(\cdot\mid y)$.
  \item $\eta$ denotes a simple reference law such as a Gaussian.
  \item $\pi$ denotes a coupling or transport plan, and $\pi_\varepsilon$ an
  entropically regularized plan.
  \item $T$ denotes a static transport map; $T_y$ denotes an observation- or
  context-conditioned map.
  \item $\rho_t$ denotes a time-indexed path of measures, and $v_t$ its velocity
  field.
  \item $a$ and $b$ denote discrete source and target marginals
  in finite-dimensional papers.
\end{itemize}

For each core paper we use the same teaching rhythm:
\begin{enumerate}[label=(\arabic*)]
  \item what problem the paper is solving,
  \item paper reproduction,
  \item explanatory reconstruction,
  \item BayesFilter bridge,
  \item and what remains unresolved.
\end{enumerate}
Whenever a step in a paper is mathematically correct but pedagogically abrupt, the
note says so and fills in the missing bridge explicitly.

\chapter{Static map learning and the OT-ICNN cautionary lesson}

\section{Why quadratic OT suggests convex potentials}

The classical quadratic-cost Monge problem is
\begin{equation}
\inf_{\pushf{T}{Q}=P} \int_{\R^d} \frac12\|x-T(x)\|^2\,dQ(x).
\label{eq:monge-problem}
\end{equation}
A first-year student should ask immediately: why do papers start talking about
convex functions instead of maps?  The answer is Brenier's theorem.  Under
regularity conditions such as finite second moments and absolute continuity of the
source law, the optimal map is the gradient of a convex potential.  The dual or
semidual formulation therefore becomes a way of learning the map indirectly through
a structured scalar object.

A semidual quadratic formulation can be written as
\begin{equation}
\W_2^2(P,Q)
=
C_{P,Q} - \inf_{f\in\mathrm{CVX}} \Bigl\{\E_{X\sim P}[f(X)] + \E_{Y\sim Q}[f^*(Y)]\Bigr\},
\label{eq:quadratic-semidual}
\end{equation}
where $f^*$ is the convex conjugate of $f$,
\begin{equation}
f^*(y)=\sup_x \{\langle x,y\rangle - f(x)\}.
\label{eq:convex-conjugate}
\end{equation}
Pedagogically, the key point is not the exact constant $C_{P,Q}$; it is the
logic:
\begin{quote}
quadratic OT turns an optimization over maps into an optimization over convex
potentials, and the map can then be recovered from the potential.
\end{quote}
That is why the OT-ICNN line exists at all.

\section{Paper reproduction: Makkuva et al. (2020)}

Makkuva et al. study quadratic-cost Monge transport with
\begin{equation}
c(x,y)=\frac12\|x-y\|^2.
\label{eq:makkuva-cost}
\end{equation}
The paper's deployed transport map is not an arbitrary regressor.  It is the
Brenier-style gradient map
\begin{equation}
T^\star = \nabla f^{\star},
\label{eq:makkuva-gradient-map}
\end{equation}
recovered from an optimized convex potential.  The paper then introduces a neural
approximation in which the Wasserstein-2 objective is approximated by a minimax
ICNN program of the form
\begin{equation}
\widetilde W_2^2(P,Q)
=
\sup_{f\in\mathrm{ICNN}}\inf_{g\in\mathrm{ICNN}} \mathcal V_{P,Q}(f,g)+C_{P,Q},
\label{eq:makkuva-minimax}
\end{equation}
which is the paper's Equation~(6) in its own notation.  The numerical method then
optimizes this minimax program by alternating updates, which is the role of the
paper's Algorithm~1.

A first implementation detail that the original paper leaves easy to miss is that
the two ICNNs are not two deployed transport maps.  The candidate map that will
ultimately be used is the gradient of the learned convex potential,
\begin{equation}
T_{\theta}(x)=\nabla f_{\theta}(x).
\label{eq:makkuva-deployed-icnn-map}
\end{equation}
The second ICNN is part of the saddle-point optimization machinery.  It helps
enforce the dual OT structure during training, but it is not the object one keeps
at deployment time.

So the operational contract is:
\begin{enumerate}[label=(\roman*)]
  \item choose two convex-network parameterizations for the minimax problem,
  \item alternate updates so that the saddle objective \eqref{eq:makkuva-minimax}
  is optimized numerically,
  \item after training, discard the optimization-side helper and deploy only the
  gradient of the learned convex potential.
\end{enumerate}
This is the paper's intended training/deployment split.

\section{Explanatory reconstruction}

The missing pedagogical bridge in the original paper is usually the same one that
confuses readers later: why should learning a convex scalar function give a good
transport map?  The answer is the chain
\begin{equation}
\text{quadratic OT}
\Longrightarrow
\text{Brenier structure}
\Longrightarrow
\text{gradient of convex potential}
\Longrightarrow
\text{ICNN parameterization}.
\label{eq:makkuva-logic-chain}
\end{equation}
So the neural net is not learning a map because ``neural nets are flexible.''  It
is learning a map because a theorem says the map should be a gradient of a convex
potential in the ideal setting, and ICNNs provide a parametric family that can
respect that structure.

A 1D version makes this less mysterious.  If $Q=\mathcal N(0,1)$ and
$P=\mathcal N(m,s^2)$, then the optimal map is
\begin{equation}
T(x)=m+sx,
\label{eq:gaussian-affine-map}
\end{equation}
and one convex potential giving this map is
\begin{equation}
\psi(x)=mx+\frac{s}{2}x^2.
\label{eq:gaussian-convex-potential}
\end{equation}
The map is the derivative:
\begin{equation}
\psi'(x)=m+sx=T(x).
\label{eq:gaussian-convex-derivative}
\end{equation}
This tiny calculation is what the Brenier story is trying to generalize.

\paragraph{What are the two ICNNs doing in practice?}
A first-year student should read the minimax program as an optimization game, not
as two separate candidate transport maps.  One ICNN parameterizes the convex
potential whose gradient is meant to become the transport map.  The other ICNN
plays the optimization-side role needed by the paper's saddle formulation.  So the
practical training loop is not ``fit a map by regression.''  It is:
\begin{align}
\text{choose candidate convex potential}
&\Longleftrightarrow
\text{improve dual witness in the saddle objective}
\notag\\
&\Longleftrightarrow
\text{alternate until the minimax objective stabilizes}.
\label{eq:makkuva-training-logic}
\end{align}
That is why the paper needs Algorithm~1 rather than a single direct supervised
loss.

\paragraph{What is optimized in practice?}
Operationally, each training iteration uses samples from the source and target
laws, evaluates the saddle objective \eqref{eq:makkuva-minimax}, performs an
update for the optimization-side ICNN, then performs an update for the candidate
transport-potential ICNN.  In a first pass implementation, the student should
therefore expect a loop of the form:
\begin{algorithm}[ht]
\caption{Makkuva-style OT-ICNN training loop}
\begin{algorithmic}[1]
\Require source samples $x\sim Q$, target samples $y\sim P$, convex networks $f_\theta,g_\omega$
\For{each iteration}
  \State sample minibatches from $P$ and $Q$
  \State evaluate the saddle objective $\mathcal V_{P,Q}(f_\theta,g_\omega)$
  \State update the optimization-side network $g_\omega$
  \State update the candidate transport-potential network $f_\theta$
\EndFor
\end{algorithmic}
\end{algorithm}
This is deliberately schematic, but it names the real executable loop rather than
leaving the reader with only the variational formula.

\paragraph{What is deployed after training?}
After training, one does not deploy the whole minimax game.  One deploys the map
obtained from the learned convex potential,
\begin{equation}
T_\theta(x)=\nabla f_\theta(x).
\label{eq:makkuva-deployment-map-again}
\end{equation}
So the deployment recipe is:
\begin{enumerate}[label=(\arabic*)]
  \item receive a new source sample $x$,
  \item evaluate the gradient of the learned convex potential,
  \item output $T_\theta(x)$ as the transported sample.
\end{enumerate}
That is the practical meaning of the Brenier-to-ICNN bridge.

\paragraph{Why this still matters for BayesFilter.}
Even though this paper is no longer the main BayesFilter lane, it remains the right
cautionary baseline for a specific reason: it teaches what a \emph{faithful}
Brenier-style direct-map implementation would actually have to contain.  The note
should therefore make clear that a faithful first pass would need
\begin{enumerate}[label=(\roman*)]
  \item a convex-potential parameterization,
  \item an alternating saddle-point training loop,
  \item and gradient-of-potential deployment.
\end{enumerate}
Without all three pieces, one has not yet reproduced the paper's algorithmic
contract.

\section{Code snapshot}

The repo's existing code-availability ledger records official public code for
OT-ICNN.  That means this paper is one of the clearest cases where students can
inspect a paper-authored implementation while reading the note.  However, the same
ledger and the current note also make clear that BayesFilter's earlier negative
result was not a faithful reproduction of that original route.

So the practical status is:
\begin{enumerate}[label=(\roman*)]
  \item official paper code exists,
  \item it is useful for faithful reproduction and for understanding the intended ICNN/Brenier training loop,
  \item but it is not yet a BayesFilter-ready filtering module without a fresh object and framework translation.
\end{enumerate}

\section{BayesFilter bridge and OT-ICNN closeout}

The current BayesFilter negative result does \emph{not} show that Makkuva et al.'s
algorithm fails.  What BayesFilter benchmarked earlier was a bridge-style surrogate,
not a faithful reproduction of the paper's full method.  The paper's central
contract is the triad
\begin{enumerate}[label=(\roman*)]
  \item convex-potential representation,
  \item minimax dual optimization,
  \item gradient-of-convex-potential deployment.
\end{enumerate}
The earlier BayesFilter bridge benchmark did not implement this full triad, so the
correct conclusion is only:
\begin{quote}
that benchmark did not beat the stronger trivial baseline.
\end{quote}
It does \emph{not} support the stronger claim that the original OT-ICNN paper is
false or ineffective.

\chapter{Conditional posterior transport: the main scientific lane}

\section{Chapter preview}

This chapter is the scientific center of gravity of the note.  The papers here are
trying to learn observation-conditioned or context-conditioned transport objects,
not merely static unconditional maps.  That is closer to what BayesFilter actually
needs.

The two hard ideas to watch are:
\begin{enumerate}[label=(\roman*)]
  \item a conditional transport problem is often a family of ordinary transport
  problems indexed by $y$,
  \item and a conditional posterior map is more demanding than a generic map,
  because it must reproduce the posterior law given the observation.
\end{enumerate}

\section{Wang et al. (2025): conditional OT for Bayesian inference}

\subsection*{What problem is the paper solving?}

Wang et al. ask whether a conditional posterior law $\pi(\cdot\mid y)$ can be
represented as the pushforward of a simple reference law by a map that depends on
$y$.  The paper then studies both a static inverse-map route and a dynamic
conditional-flow route.

For implementation purposes, the right first question is not ``what theorem is the
paper using?''  It is:
\begin{quote}
What are the inputs, what is being parameterized, what is the training objective,
and what do we evaluate at deployment time?
\end{quote}
That is how the section is organized below.

\subsection*{Paper reproduction}

The defining block-triangular map is
\begin{equation}
T(z,y)=
\begin{bmatrix}
 y \\
 g_\theta(z;y)
\end{bmatrix},
\label{eq:wang-block-triangular}
\end{equation}
which is the paper's Equation~(2.2).  Once $y$ is fixed, the conditional transport
object is just the fiberwise map $z\mapsto g_\theta(z;y)$.

If $g_\theta(\cdot;y)$ is a diffeomorphism, the conditional density follows by
change of variables:
\begin{equation}
\pushf{g_\theta(\cdot;y)}{\eta}(x)
=
\eta\bigl(g_\theta^{-1}(x;y)\bigr)
\left|\det \nabla_x g_\theta^{-1}(x;y)\right|.
\label{eq:wang-pushforward-density}
\end{equation}
The target is therefore
\begin{equation}
\pushf{g_\theta(\cdot;y)}{\eta} \approx \pi(\cdot\mid y),
\label{eq:wang-conditional-target}
\end{equation}
which is the paper's Equation~(2.4).

For the static PCP-Map route, the inverse conditional generator is fitted by the
negative log-likelihood functional
\begin{equation}
J_{\mathrm{NLL}}[g_\theta^{-1}]
=
\E_{\pi(x,y)}
\left[
\frac12\|g_\theta^{-1}(x;y)\|^2
- \log\det \nabla_x g_\theta^{-1}(x;y)
\right],
\label{eq:wang-nll}
\end{equation}
with the Brenier-style inverse parameterization
\begin{equation}
g_\theta^{-1}(x;y)=\nabla_x G_\theta(x,y).
\label{eq:wang-inverse-brenier}
\end{equation}
For the dynamic COT-Flow route, the state trajectory obeys
\begin{equation}
\dot p_t = v_\theta(t,p_t;y),
\label{eq:wang-flow-ode}
\end{equation}
and the OT-regularized objective is
\begin{equation}
\min_{v_\theta} J_{\mathrm{NLL}}[g_\theta^{-1}] + \alpha_1 P_{\mathrm{dynOT}}[v_\theta],
\label{eq:wang-flow-objective}
\end{equation}
with dynamic penalty
\begin{equation}
P_{\mathrm{dynOT}}[v_\theta]
=
\E_{\pi(x,y)}\left[\int_0^1 \frac12\|v_\theta(t,p_t;y)\|^2\,dt\right].
\label{eq:wang-dynot-penalty}
\end{equation}
The dynamic route is then related to a macroscopic density-control problem,
\begin{equation}
\min_{\rho_t,v_t}
\int \left[
-\log \rho_1(x;y)\,\pi(x,y)
+
\alpha_1\int_0^1 \frac12\|v_t(x;y)\|^2\rho_t(x;y)\,dt
\right]dx\,dy
\label{eq:wang-mfg-schematic}
\end{equation}
subject to the continuity equation
\begin{equation}
\partial_t \rho_t(x;y)+\nabla_x\cdot\bigl(\rho_t(x;y)v_t(x;y)\bigr)=0,
\qquad
\rho_0(x;y)=\eta(x).
\label{eq:wang-continuity}
\end{equation}
By Pontryagin's principle, the optimal velocity has feedback form
\begin{equation}
v_\theta(t,p_t;y) = -\frac{1}{\alpha_1}\nabla_p \Phi(t,p_t;y),
\label{eq:wang-flow-feedback}
\end{equation}
with Hamilton--Jacobi--Bellman equation
\begin{equation}
\partial_t \Phi(t,x;y)-\frac{1}{2\alpha_1}\|\nabla_x\Phi(t,x;y)\|^2=0
\label{eq:wang-hjb}
\end{equation}
and terminal condition schematically of the form
\begin{equation}
\Phi(1,x;y) = -\frac{\pi(x,y)}{\rho_1(x;y)}.
\label{eq:wang-terminal}
\end{equation}
The paper then parameterizes the value function as
\begin{equation}
\Phi_\theta(q)=\mathrm{NN}_{\theta_{\mathrm{NN}}}(q)+Q_{\theta_Q}(q),
\qquad q=(t,x;y),
\label{eq:wang-value-net}
\end{equation}
which is the role of Equation~(4.11).

\subsection*{Explanatory reconstruction}

\paragraph{Why does the block-triangular map matter?}
Equation~\eqref{eq:wang-block-triangular} is the paper's first nontrivial idea.
It means that the transport leaves the conditioning variable untouched and only
acts on the latent coordinate.  Therefore, once we freeze $Y=y$, the joint map
reduces to an ordinary transport map on the fiber over $y$:
\begin{equation}
(z,Y)\mapsto (Y,g_\theta(z;Y))
\quad\Longrightarrow\quad
z\mapsto g_\theta(z;y)
\text{ for fixed } y.
\label{eq:wang-fiberwise-reduction}
\end{equation}
So conditional transport is not a mysterious new kind of transport; it is a
family of ordinary transports indexed by the conditioning variable.

\paragraph{Why does the NLL objective have exactly that form?}
Take the reference law to be standard Gaussian,
\begin{equation}
\eta(z) = (2\pi)^{-d/2}\exp\!\left(-\frac12\|z\|^2\right).
\label{eq:gaussian-reference}
\end{equation}
Plugging \eqref{eq:gaussian-reference} into
\eqref{eq:wang-pushforward-density} gives
\begin{equation}
-\log \pushf{g_\theta(\cdot;y)}{\eta}(x)
=
\frac12\|g_\theta^{-1}(x;y)\|^2
-\log\det \nabla_x g_\theta^{-1}(x;y)
+\text{const}.
\label{eq:wang-negative-log-density}
\end{equation}
Averaging \eqref{eq:wang-negative-log-density} over samples from the target
conditional law gives exactly \eqref{eq:wang-nll}.  So the objective is not an ad
hoc design choice.  It is simply the negative log-likelihood implied by the
conditional change-of-variables formula.

\paragraph{Why the inverse map is parameterized rather than the forward map.}
The paper's static route works with $g_\theta^{-1}$ rather than directly with
$g_\theta$.  The reason is computational.  The NLL objective
\eqref{eq:wang-nll} requires the inverse map and the Jacobian determinant of that
inverse.  If the inverse is parameterized directly as a conditional Brenier map,
then the exact training objective can be evaluated without solving a numerical
inverse problem inside every gradient step.

The mathematical structure is then
\begin{equation}
\text{posterior samples }x
\xrightarrow{\ g_\theta^{-1}(\cdot;y)\ }
\text{reference samples that should look Gaussian}.
\label{eq:wang-inverse-interpretation}
\end{equation}
So the static route learns a map that pulls posterior samples backward to the
reference law.

\paragraph{What the static Brenier structure buys.}
The inverse map is constrained to be the gradient of a partially convex scalar
function,
\begin{equation}
g_\theta^{-1}(x;y)=\nabla_x G_\theta(x,y).
\label{eq:wang-static-gradient-again}
\end{equation}
This is a conditional Brenier idea: if the inverse map comes from a convex
potential in the transported variable, then it is a gradient field rather than a
free unconstrained regression map.  The paper is therefore trying to learn
``posterior samples pushed backward to look Gaussian'' in a structured way.

\paragraph{What the dynamic route is really doing.}
The dynamic route says: instead of only fitting an endpoint map, learn a whole
conditional flow whose endpoint reproduces the posterior and whose trajectories are
transport-efficient.  The penalty \eqref{eq:wang-dynot-penalty} is a kinetic-energy
penalty.  Equation \eqref{eq:wang-flow-objective} should therefore be read as a
control objective rather than as ``the NLL plus one more regularizer.''

The conceptual staircase is
\begin{align}
\text{conditional endpoint fit}
&+ \text{kinetic-energy regularization}
\notag\\
&\Longrightarrow
\text{density control problem}
\notag\\
&\Longrightarrow
\text{feedback velocity law}.
\label{eq:wang-dynamic-logic-chain}
\end{align}
So the dynamic route is not just a static map with extra smoothing.  It is a
conditional dynamic OT problem whose endpoint is the posterior law.

\paragraph{Why the HJB equation appears.}
The mean-field-game or control formulation
\eqref{eq:wang-mfg-schematic}--\eqref{eq:wang-continuity} introduces a constrained
optimization over evolving densities and velocities.  The value function
$\Phi(t,x;y)$ is the control-theoretic object that prices the best future
continuation of the transport.  Once such a value function exists, the optimal
feedback velocity is its gradient field \eqref{eq:wang-flow-feedback}, and the HJB
equation \eqref{eq:wang-hjb} is the stationarity condition of the dynamic control
problem.

This is exactly the step the original paper leaves too compressed.  The PDEs are
not arbitrary.  They are the optimality conditions of the dynamic OT control
problem.

\paragraph{What the neural network is approximating.}
The network in \eqref{eq:wang-value-net} is not directly parameterizing the final
transport map.  It parameterizes the value function from which the velocity field
is recovered.  That is a very different architectural commitment from the static
PCP-Map route.  This is the key conceptual distinction between the two halves of
the paper.

A student who gets lost in the dynamic half should stop and sort the objects by
role:
\begin{enumerate}[label=(\roman*)]
  \item $\rho_t$ is the evolving density path we are trying to fit at the population level,
  \item $v_t$ is the velocity field that moves that density path,
  \item $\Phi$ is the value-function or control object whose gradient determines the velocity,
  \item and the deployed sampler is obtained only after integrating the learned flow forward from a reference sample.
\end{enumerate}
So the learned neural object is closest to $\Phi_\theta$ (or equivalently to the
velocity field derived from it), not directly to the final posterior sample cloud.

\subsection*{Implementation manual: what Claude or a student would need}

\paragraph{Inputs.}
The static and dynamic routes both assume access to samples from the joint law
\begin{equation}
(x_i,y_i) \sim \pi(x,y),
\qquad i=1,\ldots,N,
\label{eq:wang-training-pairs}
\end{equation}
and a tractable reference law $\eta$ from which latent samples $z\sim\eta$ can be
drawn.

\paragraph{Static PCP-Map training loop.}
A practical implementation reads as follows.
\begin{algorithm}[ht]
\caption{Static PCP-Map training loop}
\begin{algorithmic}[1]
\Require minibatches $(x_i,y_i)$ from the joint law, reference law $\eta$, neural potential $G_\theta(x,y)$
\State Define $g_\theta^{-1}(x;y)=\nabla_x G_\theta(x,y)$
\State For each batch, evaluate $g_\theta^{-1}(x_i;y_i)$ and the Jacobian determinant term
\State Form the Monte Carlo estimate of $J_{\mathrm{NLL}}$ from \eqref{eq:wang-nll}
\State Backpropagate through $G_\theta$ and update $\theta$
\State At deployment, sample $z\sim\eta$ and numerically invert $g_\theta^{-1}(\cdot;y)$ or solve the corresponding forward map problem to obtain $x=g_\theta(z;y)$
\end{algorithmic}
\end{algorithm}

\paragraph{Dynamic COT-Flow training loop.}
The dynamic route instead parameterizes either the velocity field or the value
function.
\begin{algorithm}[ht]
\caption{Dynamic COT-Flow training loop}
\begin{algorithmic}[1]
\Require minibatches $(x_i,y_i)$, reference law $\eta$, value-function network $\Phi_\theta$ or velocity network $v_\theta$
\State Recover the velocity by \eqref{eq:wang-flow-feedback} if using $\Phi_\theta$
\State Simulate trajectories $p_t$ from \eqref{eq:wang-flow-ode}
\State Estimate the negative log-likelihood term and the dynamic penalty \eqref{eq:wang-dynot-penalty}
\State Backpropagate the total objective \eqref{eq:wang-flow-objective} and update parameters
\State At deployment, sample $z\sim\eta$ and integrate the ODE forward to obtain $x=g_\theta(z;y)$
\end{algorithmic}
\end{algorithm}

\paragraph{Implementation hazards.}
There are three implementation traps a student should know immediately.
\begin{enumerate}[label=(\roman*)]
  \item The static route needs stable Jacobian-determinant evaluation for the inverse map.
  \item The dynamic route needs stable ODE integration and repeated differentiation through trajectories.
  \item The two routes output different objects at deployment time: the static route centers the inverse map, while the dynamic route centers the velocity or value function.
\end{enumerate}

\paragraph{Implementation contract.}
A correct static implementation learns an \emph{inverse conditional map}
$g_\theta^{-1}(x;y)$ and deploys posterior samples only after inverting or otherwise
recovering the forward transport.  A correct dynamic implementation learns a value
function or velocity field whose endpoint flow produces the posterior sampler.  In
neither case is the raw training object itself automatically the final posterior
sample cloud.

\paragraph{Common wrong implementations.}
The easiest wrong turns are:
\begin{enumerate}[label=(\roman*)]
  \item training the inverse-map loss but deploying the raw inverse map as if it were the posterior sampler,
  \item parameterizing the forward map when the implemented loss is written for the inverse,
  \item treating $\Phi_\theta$ as the deployed artifact rather than integrating the induced flow,
  \item checking only the loss while never checking fixed-$y$ conditional pushforward quality.
\end{enumerate}
The first three change semantics, not just efficiency.

\paragraph{Minimum correct recipe.}
The shortest safe recipe is: choose the route first, fix whether the trainable object
is $g_\theta^{-1}$ or $\Phi_\theta$, train with the matching objective, and deploy
only after the required inversion or ODE integration step has been carried out.  For a
first executable build, the static inverse-map route is the safer starting point;
the dynamic route is conceptually clear here, but it still requires more numerical
decisions about trajectory integration and density-control estimation.

\paragraph{First diagnostic.}
Before any serious benchmark, freeze one observation value $y$ and run a tiny
conditional-law smoke test.  For the static route, check that pushed-back posterior
samples look Gaussian under $g_\theta^{-1}(\cdot;y)$.  For the dynamic route, check
that integrating the learned flow from reference samples at fixed $y$ actually lands
near the target conditional sample law.

\subsection*{Code snapshot}

The repo's current source-availability notes do not record a clean official codebase
for this exact paper under the title used in the note.  They do, however, record a
closely matching arXiv source for this conditional OT / Bayesian inference family.
So students should read the code status here as family-level code evidence rather
than as a clean locally endorsed implementation starting point.

Operationally, that means:
\begin{enumerate}[label=(\roman*)]
  \item this is not a theory-only family,
  \item but a student should still expect to implement the critical route-specific pieces themselves,
  \item especially because the static inverse-map route and the dynamic flow route have different deployment semantics.
\end{enumerate}

\subsection*{BayesFilter bridge}

For BayesFilter, Wang et al. are central because the paper learns exactly the kind
of object one might want in a posterior-transport lane: an
observation-conditioned sampler or inverse map.  What remains missing is not the
conditional idea, but the engineering bridge:
\begin{itemize}
  \item repeated weighted particle-cloud inputs,
  \item stable high-dimensional deployment,
  \item and downstream derivative-aware use inside BayesFilter.
\end{itemize}

\subsection*{What the paper does not settle}

The paper does not by itself establish that a learned conditional posterior map
preserves any BayesFilter-specific scalar or correction contract.  It shows how to
represent and train a conditional transport model for Bayesian inference.  The
BayesFilter-specific downstream semantics remain an implementation and validation
question.


\section{Al-Jarrah et al. (2025): nonlinear filtering with Brenier OT maps}

\subsection*{What problem is the paper solving?}

Al-Jarrah et al. ask the most BayesFilter-native question on the shelf: can a
posterior update be represented directly as an observation-conditioned transport
map?  The paper is not about generic image transport or unconditional OT.  It is
about nonlinear filtering, where one repeatedly has a predictive law and then must
condition it on a new observation.

To see the point, start from the filtering recursion.  If $\overline\pi_{t\mid t-1}$
denotes the predictive law before assimilating the new observation $Y_t$, then the
posterior law is
\begin{equation}
\pi_t = \mathcal B_{Y_t}(\overline\pi_{t\mid t-1}),
\label{eq:filtering-conditioning-operator}
\end{equation}
where $\mathcal B_y$ denotes Bayes conditioning by the observation value $y$.  The
paper's core question is therefore:
\begin{quote}
Can Bayes conditioning itself be represented as a transport operator?
\end{quote}
If the answer is yes, then posterior updating becomes a map-learning problem.

For implementation purposes, the right questions are again concrete:
\begin{enumerate}[label=(\roman*)]
  \item What samples are assumed available?
  \item What neural objects are parameterized?
  \item What population problem are they approximating?
  \item What exactly is applied to particles at deployment time?
\end{enumerate}

\subsection*{Paper reproduction}

The paper formulates the design problem directly.  For every prior or predictive
law $\mu$, find an observation-conditioned transport map $T_y$ such that
\begin{equation}
\pushf{T_y}{\mu} = \mathcal B_y(\mu),
\label{eq:aljarrah-design-problem}
\end{equation}
which is the paper's Equation~(6) in the note's notation.  This is the central
deployment equation of the whole paper.

To express this transport condition in a way that is easier to optimize, the paper
introduces independent copies $U\sim P_X$ and $\overline U\sim P_X$ together with
an observation $Y\sim P_Y$.  The bookkeeping is easier if one assigns them fixed
roles from the start:
\begin{itemize}
  \item $U$ denotes the target-side state variable whose joint law with $Y$ we want to match,
  \item $\overline U$ denotes an independent source-side sample that will actually be pushed through the learned map,
  \item $Y$ is the conditioning observation that chooses which transport map is applied.
\end{itemize}
The conditional consistency condition is then written
as
\begin{equation}
(T_y(\overline U),Y) \sim P_{U,Y},
\label{eq:aljarrah-joint-law-condition}
\end{equation}
which is the content of the paper's Equation~(8a).  Equivalently, for all suitable
test functions $F$,
\begin{equation}
\E[F(T_y(U),y)\mid y] = \E[F(U,y)\mid y].
\label{eq:aljarrah-consistency}
\end{equation}

Once this consistency condition has been fixed, the paper formulates the
conditional Monge problem
\begin{equation}
\min_{T_y} \E\bigl[c(T_y(\overline U),\overline U)\bigr]
\quad\text{subject to}\quad
(T_y(\overline U),Y)\sim P_{U,Y},
\label{eq:aljarrah-conditional-monge}
\end{equation}
with quadratic cost
\begin{equation}
c(u,u')=\frac12\|u-u'\|^2.
\label{eq:aljarrah-quadratic-cost}
\end{equation}
The paper then rewrites the problem by duality as a conditional max--min game,
\begin{equation}
\sup_f\inf_{T_y} J(f,T_y),
\label{eq:aljarrah-minimax-abstract}
\end{equation}
and more explicitly as
\begin{equation}
J(f,T_y)
=
\E[f(U',y)]
+
\E\bigl[c(T_y(U),U)-f(T_y(U),y)\bigr].
\label{eq:aljarrah-max-min}
\end{equation}
The empirical objective is then
\begin{equation}
\widehat J_N(f,T_y)
=
\frac1N\sum_{i=1}^{N} f(U_i',y)
+
\frac1N\sum_{i=1}^{N}\bigl[c(T_y(U_i),U_i)-f(T_y(U_i),y)\bigr],
\label{eq:aljarrah-empirical-objective}
\end{equation}
which is the finite-sample problem the neural networks actually optimize.

Finally, the deployed posterior-update operator is
\begin{equation}
\Pi_y(\mu)=\pushf{T_y}{\mu},
\label{eq:aljarrah-posterior-map}
\end{equation}
which is the BayesFilter-relevant object.

\subsection*{Explanatory reconstruction}

\paragraph{Why is the transport problem conditional rather than ordinary?}
A first reading of the paper can be baffling because the transport target is not a
single fixed measure.  It is the posterior law, which changes with the
observation.  Equation~\eqref{eq:aljarrah-design-problem} is therefore not saying
``learn one map from prior to posterior.''  It is saying:
\begin{quote}
for each observation value $y$, learn a different map $T_y$ whose pushforward is
exactly the conditioned law.
\end{quote}
So the learned object is a family of maps indexed by $y$, not one static map.

\paragraph{Why does the test-function identity characterize the posterior law?}
Suppose \eqref{eq:aljarrah-consistency} holds for all bounded measurable test
functions $F$.  Then, by the defining property of equality in law,
\begin{equation}
\E[F(X_1)\mid y]=\E[F(X_2)\mid y]\ \forall F
\quad\Longrightarrow\quad
\mathcal L(X_1\mid y)=\mathcal L(X_2\mid y).
\label{eq:test-function-law-identity}
\end{equation}
This is the crucial bridge many readers need.  The paper is not only matching a
few posterior moments.  It is matching the full conditional law via a test-function
identity.

A useful mental model is the following.  If no test function can distinguish
$T_y(U)$ from the target posterior variable, then the two have the same posterior
law.  So the conditional consistency condition is not decorative; it is the core
correctness condition.

\paragraph{Why does the Monge problem appear?}
Once the law-matching condition is fixed, there are typically many maps that send
$P_X$ to $P_{X\mid Y=y}$.  The paper therefore chooses among them by minimizing a
quadratic transport cost.  This is exactly the role of
\eqref{eq:aljarrah-conditional-monge}: among all observation-conditioned maps that
produce the right conditional law, pick the one that moves mass least under the
quadratic cost.  So the conditioning problem is being turned into a conditional OT
selection principle.

\paragraph{What are the critic and the map doing in the max--min game?}
The max--min formulation can be hard to read on first sight.  Here is the slow
interpretation.

- The critic $f$ tries to score whether transported points look like they come from
  the target joint law.
- The map $T_y$ tries to move source points into locations that the critic accepts,
  while paying the transport cost.

Thus
\begin{equation}
\text{critic certifies target agreement}
\quad\Longleftrightarrow\quad
\text{map pays transport cost to achieve it}.
\label{eq:critic-map-division}
\end{equation}
The first expectation in \eqref{eq:aljarrah-max-min} rewards the critic on true
posterior-side samples, while the second expectation penalizes the map by its OT
cost and subtracts the critic score at the transported location.  This is the
conditional OT analogue of the usual primal-dual game, but now with the observation
variable carried along.

\paragraph{Why is the empirical objective the right sample version?}
The empirical objective \eqref{eq:aljarrah-empirical-objective} is simply the
sample-average version of the population objective \eqref{eq:aljarrah-max-min}.
This matters because the paper is designed to work in a likelihood-free or
simulation-based setting: one only needs samples from the joint or conditional
observation mechanism, not a closed-form likelihood expression.  So the neural
training problem is not detached from the theory; it is the Monte Carlo version of
that theory.

\paragraph{What hidden assumption matters most?}
The hardest hidden step is the role of convexity.  The theory later needs
conditions like strong convexity of the function
\begin{equation}
u \mapsto \frac12\|u\|^2 - f(u,y)
\label{eq:aljarrah-strong-convexity-bridge}
\end{equation}
with respect to the transported variable.  Why does this matter?  Because that is
what allows one to control the difference between an approximate learned pair
$(f,T_y)$ and the ideal optimal pair.  Without such structure, the max--min problem
might still be optimized numerically, but the map-quality guarantees become much
weaker.

\paragraph{What the paper leaves implicit.}
The original paper moves quickly from Bayes-law motivation to the conditional OT
game.  The missing bridge supplied here is:
\begin{align}
\text{Bayes conditioning}
&\Longrightarrow
\text{conditional law identity}
\notag\\
&\Longrightarrow
\text{conditional Monge selection principle}
\notag\\
&\Longrightarrow
\text{dual or max--min learning problem}.
\label{eq:aljarrah-logic-chain}
\end{align}
Once this ladder is explicit, the whole paper becomes much easier to follow.

\subsection*{Implementation manual: what Claude or a student would need}

\paragraph{Inputs.}
A practical implementation needs predictive or prior particles together with
observation samples.  In the paper's finite-sample construction, one needs
triples of the form
\begin{equation}
(U_i,\overline U_i,Y_i),
\qquad
U_i,\overline U_i \sim P_X,
\quad
Y_i\sim P_Y
\text{ or }Y_i\sim h(\cdot\mid U_i),
\label{eq:aljarrah-training-triples}
\end{equation}
where $\overline U_i$ is an independent source-side copy and $Y_i$ is the
conditioning observation paired with the target-side sample.

\paragraph{Parameterization.}
The paper uses two learnable objects:
\begin{enumerate}[label=(\roman*)]
  \item a critic $f_\phi(u,y)$,
  \item and an observation-conditioned transport map $T_\theta(u,y)$,
\end{enumerate}
with the map represented by a residual network and the critic by a scalar-valued
network.

\paragraph{Population objective.}
The ideal optimization target is the max--min program
\eqref{eq:aljarrah-minimax-abstract}--\eqref{eq:aljarrah-max-min}.  This is what
should be written down before coding anything else, because it tells us exactly
which expectations must be estimated and which variables are being optimized.

\paragraph{Empirical objective.}
The code should implement the Monte Carlo objective
\eqref{eq:aljarrah-empirical-objective}.  In practice this means:
\begin{enumerate}[label=(\arabic*)]
  \item draw or assemble a batch of predictive-state samples $U_i$,
  \item form an independent batch $\overline U_i$ or a shuffled copy,
  \item sample or simulate the corresponding observations $Y_i$,
  \item evaluate the critic on target-side samples,
  \item evaluate the transported samples $T_\theta(\overline U_i,Y_i)$,
  \item accumulate the transport-cost and critic terms,
  \item update the neural parameters.
\end{enumerate}
If the symbols start to blur together, keep the intended semantics in one sentence:
\begin{quote}
$U_i$ is the law we want to match, $\overline U_i$ is the independent source sample
we are allowed to move, and $Y_i$ tells the network which posterior-update map to
use.
\end{quote}

\paragraph{Training loop.}
A minimal implementation loop is:
\begin{algorithm}[ht]
\caption{Conditional Brenier OT filter-map training loop}
\begin{algorithmic}[1]
\Require batches $(U_i,\overline U_i,Y_i)$, critic network $f_\phi$, transport network $T_\theta$, quadratic cost $c$
\State Compute transported samples $\widehat U_i=T_\theta(\overline U_i,Y_i)$
\State Evaluate the empirical objective $\widehat J_N(f_\phi,T_\theta)$ from \eqref{eq:aljarrah-empirical-objective}
\State Update $\phi$ to increase the critic objective
\State Update $\theta$ to decrease the transport objective
\State Repeat until the min--max game stabilizes
\end{algorithmic}
\end{algorithm}
The paper's actual implementation uses stochastic optimization and residual-network
parameterizations, but this is the core mathematical loop.

\paragraph{Deployment loop.}
At deployment time, the algorithmic object is not the critic.  It is the learned
posterior-update operator
\begin{equation}
\mu \longmapsto \Pi_y(\mu)=\pushf{T_y}{\mu}.
\label{eq:aljarrah-deployment-loop}
\end{equation}
So a BayesFilter-style use is straightforward in form:
\begin{enumerate}[label=(\arabic*)]
  \item start with predictive particles $u_i$ approximating $\mu$,
  \item observe $y$,
  \item push every predictive particle through $T_y$,
  \item treat the transported cloud as the posterior particle approximation.
\end{enumerate}

\paragraph{Implementation hazards.}
A reader trying to code the paper should be warned about four difficulties.
\begin{enumerate}[label=(\roman*)]
  \item The correctness of the learned map depends on the conditional consistency
  structure, not only on numerical loss decrease.
  \item The theory relies on convexity or strong-convexity structure that the raw
  network parameterization may not enforce automatically.
  \item The map is conditioned on $y$, so batching and sample pairing must preserve
  the conditioning semantics.
  \item The finite-sample construction may be easier to state than to stabilize in
  high dimensions.
\end{enumerate}

\paragraph{Implementation contract.}
The learned object is the family of observation-conditioned maps $T_y$; the optimized
objects are the map and critic in the max--min game; the critic is auxiliary only; and
the deployed object is the posterior-update operator $\Pi_y(\mu)=\pushf{T_y}{\mu}$.
A correct implementation must preserve the conditional-law identity, not only obtain a
small adversarial or primal-dual loss.

\paragraph{Common wrong implementations.}
The most likely wrong turns are:
\begin{enumerate}[label=(\roman*)]
  \item confusing $U$, $\overline U$, and $Y$ so the batch semantics are wrong,
  \item learning one static map instead of a family $T_y$ indexed by the observation,
  \item deploying the critic or critic score instead of deploying only the transported cloud,
  \item treating lower min--max loss as sufficient without checking the conditional-law contract.
\end{enumerate}
The first three are semantic failures, not just tuning mistakes.

\paragraph{Minimum correct recipe.}
The shortest safe recipe is: build correctly paired triples $(U_i,\overline U_i,Y_i)$,
train the critic and conditional map against the empirical max--min objective, discard
the critic at deployment, and use only $T_y$ to move predictive particles at fixed
observation value $y$.

\paragraph{First diagnostic.}
Before any large run, freeze one observation value $y$ and verify that transported
particles are hard to distinguish from the target conditional law at that same $y$.
If batch shuffling or conditioning semantics are wrong, this fixed-$y$ check usually
fails immediately even when the training loss decreases.

\subsection*{Code snapshot}

The repo's existing ledger does not record a confirmed public implementation source
for this filtering-native Brenier OT paper.  Students should therefore interpret the
status as conceptually strong but practically paper-first: the note itself is doing
much of the implementation-guidance work.

The practical interpretation is:
\begin{enumerate}[label=(\roman*)]
  \item there is no confirmed local code-backed shortcut,
  \item the implementation contract in this section matters more than any presumed external repo,
  \item and a first BayesFilter implementation should be treated as a fresh reconstruction.
\end{enumerate}

\subsection*{BayesFilter bridge}

This paper is extremely important for BayesFilter because the deployed object
\eqref{eq:aljarrah-posterior-map} is already a posterior-update operator.  That is
much closer to BayesFilter's real need than a generic static OT map.  The remaining
gap is that the learned conditional map still needs to be translated into a stable,
reusable, derivative-aware filtering block.

\subsection*{What the paper does not settle}

The paper does not prove that its learned map preserves all downstream BayesFilter
semantics beyond the conditional transport contract it studies.  It is therefore a
strong scientific lane, but not yet an automatic drop-in module.


\section{Hosseini, Hsu, and Taghvaei (2025): conditional OT on function spaces}

\subsection*{What problem is the paper solving?}

Hosseini et al. ask the theorem-level version of the same conditional question:
can conditional transport be formulated directly on function spaces, not only in
finite-dimensional Euclidean coordinates?  This makes the paper conceptually very
important, but also harder to read.

A first-year student should notice immediately what is unusual here.  In Wang, the
main burden is conditional change of variables inside a finite-dimensional neural
map.  Here the burden is different: the paper first asks what the \emph{ideal
conditional transport object} should be in infinite-dimensional settings, and only
later discusses approximation.

If the implementation-ready question is asked immediately, the core difficulty
becomes clear:
\begin{quote}
What is the algorithm supposed to approximate if the paper starts by formulating a
theorem-level transport object on function spaces?
\end{quote}
The section below is written to answer that exact question.

\subsection*{Paper reproduction}

The starting point is the paper's conditional transport problem.  In the note's
notation, we want a conditional map $T_y^\star$ such that
\begin{equation}
\pushf{T_y^\star}{\eta} = \pi(\cdot\mid y).
\label{eq:hosseini-pushforward}
\end{equation}
This is the paper's Problem~1.1 in house notation.  The significance is that the
map does not target one fixed distribution.  It targets a whole family of
conditionals.

The paper then places this inside a block-triangular map.  In the notation used
throughout the note,
\begin{equation}
T(y,u) = (y,T_y(u)).
\label{eq:hosseini-block-triangular}
\end{equation}
This is the finite-dimensional shadow of the paper's Definition~1.2.  The
triangular structure means that the conditioning variable is carried through by the
identity map while only the latent or state coordinate is transported.

To formulate the Bayesian inverse problem, the paper considers a joint law of data
and unknowns.  In house notation, the posterior law appears schematically as
\begin{equation}
\pi(du\mid y) \propto L(y\mid u)\,\mu_0(du),
\label{eq:hosseini-bip-posterior}
\end{equation}
which is the conceptual content of the paper's Equations~(1.3)--(1.4).  The key
point is that conditional OT is being used to characterize this posterior law as a
transport image of a reference law, not merely as a Radon--Nikodym derivative.

The conditional Monge problem then becomes
\begin{equation}
\inf_{T_y}
\int c(u,T_y(u))\,d\eta_y(u)
\quad\text{subject to}\quad
\pushf{T_y}{\eta_y} = \pi(\cdot\mid y)
\text{ for a.e. } y,
\label{eq:hosseini-conditional-monge}
\end{equation}
which is the note's translation of Problem~1.3.  The conditional Kantorovich
relaxation then replaces the map by a constrained coupling that keeps the data
coordinate fixed; this is the content of Problem~1.4.

Later, in the explicit Bayesian-inverse-problem section, the paper specializes the
reference law to an independence coupling between the data marginal and the prior,
so that the conditional Monge problem becomes, schematically,
\begin{equation}
\inf_{T}
\int \|v-T_y(v)\|_{\mathcal U}^{p}\,d\eta(y,v)
\quad\text{subject to}\quad
\pushf{T}{\eta}=\nu,
\qquad
T(y,v)=(y,T_y(v)),
\label{eq:hosseini-bip-monge}
\end{equation}
which is the role of the paper's Equation~(4.2).  The point of this specialization
is to show that the same conditional OT machinery can represent a prior-to-posterior
transport family for inverse problems.

Finally, the paper proposes two approximation routes later in the manuscript:
Algorithm~5.1 (plugin conditional OT) and Algorithm~5.2 (WaMGAN).  The theory and
the approximation are deliberately not the same object.  That separation is one of
the paper's main virtues.

\subsection*{Explanatory reconstruction}

\paragraph{First do the finite-dimensional version.}
The paper becomes much easier if one first ignores function spaces and imagines
ordinary Euclidean variables.  Fix an observation value $y$.  Then the paper is
trying to find a map that sends a simple reference law to the posterior law for
that particular observation:
\begin{equation}
\eta_y \xrightarrow{\ T_y\ } \pi(\cdot\mid y).
\label{eq:hosseini-fiber-picture}
\end{equation}
Now vary $y$.  Instead of one transport problem, you now have a whole family of
transport problems, one for each observation fiber.  This is the real meaning of
conditional OT in the paper.

\paragraph{Why the triangular structure matters.}
A student might ask: why not just learn a generic map on the pair $(y,u)$?  The
triangular structure in \eqref{eq:hosseini-block-triangular} answers this.  If the
map is of the form
\begin{equation}
(y,u)\mapsto (y,T_y(u)),
\label{eq:triangular-explicit}
\end{equation}
then the data variable is not transported.  It only indexes which transport map is
used on the latent variable.  This is the mathematical mechanism that turns a joint
transport problem into a conditional one.

Without triangularity, the map could also move the conditioning variable itself,
and then the interpretation ``posterior law given the same observation $y$'' would
break down.  So the triangular restriction is not cosmetic.  It is what makes the
problem genuinely conditional.

\paragraph{What does ``for almost every $y$'' mean?}
Equation \eqref{eq:hosseini-conditional-monge} requires
\begin{equation}
\pushf{T_y}{\eta_y} = \pi(\cdot\mid y)
\quad\text{for a.e. }y.
\label{eq:hosseini-ae-condition}
\end{equation}
This means the transport condition can fail only on a set of observations of
probability zero.  Operationally, it is the correct conditional analogue of saying
that an ordinary pushforward identity holds almost surely.  The paper is therefore
not demanding a uniform deterministic identity for every logically possible $y$,
but for all observations that matter under the underlying measure.

\paragraph{Why the conditional Monge problem is the right object.}
Once the target posterior family $\pi(\cdot\mid y)$ has been fixed, there may be
many maps that realize it.  The conditional Monge problem
\eqref{eq:hosseini-conditional-monge} picks a preferred one by minimizing the
transport cost for each fiber.  So the logic is:
\begin{align}
\text{conditional law matching}
&\Longrightarrow
\text{many admissible transport maps}
\notag\\
&\Longrightarrow
\text{choose the least-cost one fiberwise}.
\label{eq:hosseini-selection-logic}
\end{align}
That is the hidden selection principle behind the paper.

\paragraph{Why the Kantorovich relaxation is needed.}
A first-year student should ask why the paper introduces a conditional coupling at
all.  The answer is the same as in ordinary OT: deterministic maps can be hard to
prove exist, hard to optimize directly, or too rigid in approximation arguments.
The constrained conditional coupling problem is the natural relaxation of the
conditional Monge problem.  The novelty is not the relaxation itself; it is that
the coupling is constrained so the data variable is still deterministically matched
by the identity map.

\paragraph{What is new in the function-space setting?}
The finite-dimensional logic above is already enough to understand the architecture
of the argument.  The function-space leap adds three new difficulties:
\begin{enumerate}[label=(\roman*)]
  \item the ambient spaces are separable Hilbert spaces,
  \item the cost and regularity assumptions must now be stated in infinite-dimensional form,
  \item and the eventual algorithm must pass through finite truncation or approximation.
\end{enumerate}
So the function-space paper is not conceptually new in its conditional logic.  It
is mathematically harder because it proves that the same logic survives in a much
richer setting.

\paragraph{Why is this paper harder than Wang?}
Wang begins with a neural parameterization and then fits it by likelihood and OT
regularization.  Hosseini et al. do the reverse.  They first define the ideal
conditional OT object and only later discuss how to approximate it.  The paper's
real contribution is therefore the separation
\begin{equation}
\text{ideal conditional transport object}
\quad\Longrightarrow\quad
\text{approximation route for that object}.
\label{eq:hosseini-theory-to-approximation}
\end{equation}
That separation is what makes the paper valuable, but it is also what makes the
reading harder.  The student is being asked to keep the theorem-level and
algorithm-level objects separate.

\paragraph{What the plugin and WaMGAN algorithms are really doing.}
The plugin route says: estimate the conditional target structure first, then solve
conditional OT numerically.  The WaMGAN route says: parameterize a conditional
generator and train it adversarially so it behaves like the conditional OT object.
So the two algorithms answer different practical questions:
\begin{align}
\text{plugin route:}
&\quad
\text{estimate target then transport};
\notag\\
\text{WaMGAN route:}
&\quad
\text{learn a generator that emulates transport}.
\label{eq:hosseini-two-algorithms}
\end{align}
This distinction is easy to miss in a fast reading of the paper, but it is very
important pedagogically.

The next question a student should ask is: where does the finite implementation
burden enter?  The answer is that the theorem-level object lives on an infinite-
dimensional space, while the code must choose a finite representation, a finite
family of conditioning values, and a finite numerical surrogate for the conditional
OT solve.  So the real bridge is not
\begin{quote}
``the theorem immediately gives an algorithm,''
\end{quote}
but rather
\begin{align}
\text{ideal conditional OT object on function space}
&\Longrightarrow
\text{choose a finite representation}
\notag\\
&\Longrightarrow
\text{approximate the conditional target family}
\notag\\
&\Longrightarrow
\text{solve OT fiberwise or train a surrogate generator}.
\label{eq:hosseini-finite-bridge}
\end{align}
That ladder is the practical content hidden between the theorem section and the two
algorithms.

\paragraph{What the paper leaves implicit.}
For new readers, the hardest gap is not the notation; it is the leap from a
finite-dimensional conditional map intuition to a function-space theorem.  The
correct reading order is:
\begin{align}
\text{fiberwise conditional OT in finite dimensions}
&\Longrightarrow
\text{triangular conditional transport}
\notag\\
&\Longrightarrow
\text{function-space generalization}.
\label{eq:hosseini-reading-order}
\end{align}
If one tries to read the function-space theorem first, the paper feels mysterious.
If one reads the finite-dimensional logic first, the abstraction becomes much more
reasonable.

A useful concrete mental model is a finite truncation of an unknown function.
Imagine that the hidden state is represented approximately by its first few basis
coefficients,
\begin{equation}
u \approx (u_1,\ldots,u_m) \in \R^m.
\label{eq:hosseini-finite-truncation-idea}
\end{equation}
Then the theorem-level object $T_y^\star$ says: for each observation value $y$,
there should be a map on this coefficient vector that pushes a simple reference law
to the posterior law of those coefficients.  The plugin route would try to solve
that finite-dimensional conditional OT problem more directly; the WaMGAN route
would try to learn a generator that emulates the same conditional pushforward.
This finite-truncation picture is not the theorem itself, but it is the right first
mental model for implementation.

\subsection*{Implementation manual: what Claude or a student would need}

\paragraph{Inputs.}
The theorem object assumes a joint measure of observations and unknowns
\begin{equation}
\nu \in \calP(\mathcal Y\times \mathcal U)
\label{eq:hosseini-joint-measure}
\end{equation}
and a reference law of the form
\begin{equation}
\eta = \nu_{\mathcal Y}\otimes \eta_{\mathcal V}
\quad\text{or, in the inverse-problem specialization,}\quad
\eta = \nu_{\mathcal Y}\otimes \mu.
\label{eq:hosseini-reference-law}
\end{equation}
So an implementer must know whether the reference variable is a synthetic latent
law $\eta_{\mathcal V}$ or the prior itself.

\paragraph{Theorem object versus algorithmic object.}
The theorem object is the ideal conditional map $T_y^\star$ satisfying
\eqref{eq:hosseini-pushforward}.  The algorithmic object is not the same thing.
It is either:
\begin{enumerate}[label=(\roman*)]
  \item a plugin numerical approximation to that map, or
  \item a trained conditional generator meant to emulate it.
\end{enumerate}
A correct implementation note must keep those two levels separate.

\paragraph{Plugin route.}
The plugin route is the more literal implementation of the theorem story.  Its
steps are approximately:
\begin{algorithm}[ht]
\caption{Plugin conditional OT pipeline}
\begin{algorithmic}[1]
\Require samples $(y_j,u_j)\sim \nu$, reference law $\eta=\nu_{\mathcal Y}\otimes\eta_{\mathcal V}$ or $\nu_{\mathcal Y}\otimes\mu$
\State Choose a finite state representation and a finite representation of the conditioning variable
\State Estimate the conditional target family $\pi(\cdot\mid y)$ or its finite projection on that representation
\State For each conditioning value or each local approximation point $y$, solve the corresponding conditional OT problem
\State Extract the $\mathcal U$-component of the resulting triangular map
\State Use interpolation or nearest-neighbor style selection to evaluate the map at new conditioning values
\end{algorithmic}
\end{algorithm}
This route is closest to the theorem, but can be computationally heavy.

\paragraph{WaMGAN route.}
The WaMGAN route is the more neural approximation route.  Its rough logic is:
\begin{algorithm}[ht]
\caption{Conditional WaMGAN-style transport training}
\begin{algorithmic}[1]
\Require samples $(y_j,u_j)\sim\nu$, chosen conditional generator family, regularization parameters
\State Choose the finite representation in which the conditional generator and critic will operate
\State Sample conditioning values and latent/reference variables
\State Push the latent variables through the conditional generator to produce synthetic conditional samples
\State Compare synthetic and target conditional samples by the paper's adversarial or Wasserstein objective
\State Update the generator (and critic, if present) by stochastic optimization
\State At deployment, sample the reference variable and evaluate the conditional generator at a new observation value $y$
\end{algorithmic}
\end{algorithm}
This route is farther from the theorem object, but easier to amortize.

\paragraph{Deployment.}
After training, the implementable object should look like
\begin{equation}
(y,v) \longmapsto T_y(v),
\qquad
v\sim \eta_{\mathcal V} \text{ or } v\sim \mu,
\label{eq:hosseini-deployment-object}
\end{equation}
so that conditional samples are obtained by fixing $y$ and pushing forward the
reference or prior law.

\paragraph{Implementation hazards.}
The major hazards are:
\begin{enumerate}[label=(\roman*)]
  \item the theorem object lives on function spaces, but the implementation must choose a finite representation,
  \item conditional OT must be evaluated or approximated fiberwise in $y$,
  \item plugin methods can become expensive if many conditioning values are needed,
  \item neural approximation methods may obscure the distinction between the theorem object and the trained surrogate,
  \item the paper's most natural inverse-problem interpretation uses a prior-based reference law, which is not the same as the Gaussian-reference logic in Wang.
\end{enumerate}

\paragraph{Implementation contract.}
The theorem object is the ideal conditional map $T_y^\star$; the algorithmic object is
either a plugin finite approximation of that map or a trained generator that emulates
it; and the deployed object is whatever finite surrogate finally sends the chosen
reference representation to conditional samples at fixed $y$.  A correct implementation
must keep those three levels separate.

\paragraph{Common wrong implementations.}
The most likely wrong turns are:
\begin{enumerate}[label=(\roman*)]
  \item treating the theorem object as if it already specified a finite algorithm,
  \item choosing a neural generator without recording the finite state representation it is supposed to emulate,
  \item pooling over conditioning values instead of respecting the fiberwise nature of the conditional problem,
  \item silently switching the reference law from prior-based to Gaussian or vice versa.
\end{enumerate}
The first and fourth mistakes change the mathematical contract, not only the code path.

\paragraph{Minimum correct recipe.}
The shortest safe recipe is: choose the finite state representation first, choose the
reference law second, decide whether the route is plugin or surrogate third, and only
then code the transport solver or generator for fixed-$y$ conditional problems.  If a
student wants one recommended first surrogate, the safest first build is a finite-
dimensional plugin route on a low-dimensional coefficient truncation before attempting
a fully learned generator.

\paragraph{First diagnostic.}
Before any broad experiment, test one finite truncation and one fixed observation value
$y$.  Check whether the implemented surrogate actually pushes the chosen finite
reference representation toward the intended conditional target for that same $y$.
This is the smallest theorem-object versus surrogate-object sanity check.

\subsection*{Code snapshot}

The repo's existing availability research does not record a confirmed official
codebase for this function-space conditional OT paper.  Students should therefore
interpret the paper as theorem-backed but effectively implementation-first in our
current notes.

In practice, the code status means:
\begin{enumerate}[label=(\roman*)]
  \item the note's theorem-to-surrogate bridge is currently more reusable than any confirmed external software source,
  \item a first implementation still requires explicit finite representation choices,
  \item and code reuse is less decisive here than avoiding theorem-object versus surrogate confusion.
\end{enumerate}

\subsection*{BayesFilter bridge}

This paper matters most if BayesFilter's long-run ambition is truly
function-space or operator-valued posterior transport.  The paper's weakness is not
lack of rigor.  It is that a large finite-dimensional bridge still has to be built
before one gets an implementable filtering routine.

\subsection*{What the paper does not settle}

The paper does not automatically give a practical weighted-particle BayesFilter
update.  It gives the strongest mathematical formulation of what a conditional
transport object should be in a richer setting.


\section{Bunne, Krause, and Cuturi (2022): a shorter conditional comparison}

Bunne et al. are useful as a contrast paper.  The learned object is a
context-conditioned Monge map family,
\begin{equation}
T_\theta(y)(u)=\nabla_u F_\theta(u,y),
\label{eq:bunne-conditional-map}
\end{equation}
with shared-task objective
\begin{equation}
\max_\theta \sum_{i=1}^{m} J_i\bigl(T_\theta(y_i)\bigr).
\label{eq:bunne-task-aggregate}
\end{equation}
This teaches a useful idea: conditional transport across a family of contexts can
be learned jointly.  But the paper is less central than Wang or Al-Jarrah because
its deployed object is still mainly a context-conditioned sample pushforward,
rather than a filtering-native posterior update or a retained teacher.

The right way to read this section is therefore not
\begin{quote}
``here is another full implementation lane to copy first.''
\end{quote}
It is instead:
\begin{quote}
``here is a shorter comparison paper that teaches the multi-context learning idea,
but does not carry the full filtering-native or retained-teacher burden.''
\end{quote}

Operationally, the learned object is a family of maps indexed by context $y$; the
objective says those context-specific tasks are trained jointly rather than one by
one.  The section's teaching value is the contrast:
\begin{enumerate}[label=(\roman*)]
  \item Wang emphasizes conditional density modeling and change of variables,
  \item Al-Jarrah emphasizes posterior-update semantics,
  \item Bunne emphasizes joint learning across many related contexts.
\end{enumerate}
So Bunne belongs here as a comparison marker, not as the main implementation anchor
for BayesFilter.

\section{What this chapter established}

The conditional lane contains the most scientifically central papers for
BayesFilter.  A first-year reader should now know:
\begin{enumerate}[label=(\arabic*)]
  \item Wang teaches how a conditional posterior map can be turned into a density
  model by change of variables.
  \item Al-Jarrah teaches how a posterior-update map can be characterized by a
  conditional law identity.
  \item Hosseini teaches the theorem-level conditional OT object that those more
  practical papers are approximating.
\end{enumerate}

\chapter{Retained-teacher acceleration: keep the solve, learn the warm start}

\section{Why this lane is conservative}

The retained-teacher lane is the right lane when you do not want the learned model
to replace OT outright.  Instead, you want the learned model to help solve the same
teacher problem faster.

To make this concrete, suppose a discrete entropic OT problem is posed over source
and target marginals $a,b$ and cost matrix $C$.  The regularized plan is
\begin{equation}
\pi_\varepsilon
\in
\argmin_{\pi\in\Pi(a,b)}
\langle C,\pi\rangle + \varepsilon \KL(\pi\,\|\,a\otimes b).
\label{eq:entropic-teacher}
\end{equation}
A retained-teacher learner does \emph{not} replace \eqref{eq:entropic-teacher}.
It predicts latent quantities that make solving \eqref{eq:entropic-teacher} easier.

\section{A tiny discrete warm-start example}

Suppose the cost matrix is $2\times 2$ and the teacher solve uses Sinkhorn-style
updates.  Then the raw optimization variable is the plan $\pi$, but the practical
solver often works through dual or scaling variables.  The idea of retained-teacher
learning is therefore:
\begin{equation}
(a,b,C)
\xrightarrow{\text{neural prediction}}
\widehat h_\theta
\xrightarrow{\text{teacher correction}}
\pi_\varepsilon.
\label{eq:warm-start-tiny-example}
\end{equation}
This is the right toy picture to keep in mind when reading Meta OT and UNOT.

\subsection*{Explanatory reconstruction}

\paragraph{What is a retained-teacher warm start actually doing?}
The most useful way to read retained-teacher papers is not by the slogan
``predict the dual.''  Start instead from the teacher solve
\eqref{eq:entropic-teacher}.  A direct solve computes the full corrected transport
plan from scratch.  A retained-teacher learner asks whether one can predict a
structured latent state from which the same teacher solve converges faster.  The
chain
\begin{equation}
(a,b,C)
\Longrightarrow
\widehat h_\theta
\Longrightarrow
\text{corrected solve for }\pi_\varepsilon
\label{eq:retained-teacher-chain}
\end{equation}
is therefore the real conceptual object, and \eqref{eq:metaot-entropic-dual} and
\eqref{eq:unot-dual-abstract} are two different mathematical ways to realize that
same retained-teacher idea.

\section{Meta OT (Amos et al., 2023)}

\subsection*{What problem is the paper solving?}

Meta OT asks whether many related OT problems can be solved faster by amortizing
the dual solution itself.  This is the cleanest retained-teacher paper on the
shelf, because it is explicit that the point is not to replace the teacher solve,
but to learn how to initialize it from the context of previous problems.

For implementation purposes, the key questions are:
\begin{enumerate}[label=(\roman*)]
  \item what teacher problem is fixed,
  \item what latent solver state is predicted,
  \item how the prediction is trained,
  \item and how a corrective solve refines that prediction.
\end{enumerate}

\subsection*{Paper reproduction}

In the discrete entropic setting, the teacher dual is
\begin{equation}
f^*,g^* \in
\argmax_{f,g}
\left\{
\langle f,a\rangle + \langle g,b\rangle
-\varepsilon \exp(f/\varepsilon)^{\top} K \exp(g/\varepsilon)
\right\},
\label{eq:metaot-entropic-dual}
\end{equation}
where $K_{ij}=\exp(-C_{ij}/\varepsilon)$.  Once the duals are known, the primal
plan is recovered by
\begin{equation}
\pi^*_{ij} = \exp(f_i^*/\varepsilon)K_{ij}\exp(g_j^*/\varepsilon),
\label{eq:metaot-primal-recovery}
\end{equation}
which is the paper's Equation~(6).

The Sinkhorn teacher itself can be written as the log-space coordinate ascent
updates
\begin{equation}
g \leftarrow \varepsilon \log b - \varepsilon \log\bigl(K^{\top}\exp(f/\varepsilon)\bigr),
\qquad
f \leftarrow \varepsilon \log a - \varepsilon \log\bigl(K\exp(g/\varepsilon)\bigr),
\label{eq:metaot-sinkhorn-updates}
\end{equation}
which is the substance of the paper's Algorithm~1.

The paper then exploits the fact that one dual potential determines the other via
\begin{equation}
g(f;b,C) = \varepsilon \log b - \varepsilon \log\bigl(K^{\top}\exp(f/\varepsilon)\bigr),
\label{eq:metaot-dual-map}
\end{equation}
so that one may amortize only a single dual variable.  This yields the objective
\begin{equation}
f^*(a,b,C,\varepsilon)
\in
\argmin_{f\in\R^n} J(f;a,b,C),
\label{eq:metaot-single-dual-objective}
\end{equation}
which is the paper's Equation~(16), and the amortization model
\begin{equation}
(a,b,C) \longmapsto \widehat f_\theta(a,b,C).
\label{eq:metaot-predicted-discrete-dual}
\end{equation}
The learning objective is then
\begin{equation}
\min_\theta \E_{(a,b,C)\sim\mathcal D} J\bigl(\widehat f_\theta(a,b,C);a,b,C\bigr),
\label{eq:metaot-amortization-loss}
\end{equation}
which is the paper's Equation~(17).

On the continuous side, the teacher problem is quadratic-cost Wasserstein-2
transport between measures $\mu$ and $\nu$.  The transport map is recovered from a
convex dual potential,
\begin{equation}
T^*(x)=\nabla \psi^*(x),
\label{eq:metaot-continuous-map}
\end{equation}
corresponding to the paper's Equation~(12), while the neural OT teacher is taken
from a W2GN-style objective.  The amortization model then predicts parameters of a
convex potential or dual object,
\begin{equation}
(\mu,\nu) \longmapsto \widehat\psi_\theta,
\label{eq:metaot-predicted-continuous-dual}
\end{equation}
and is trained by an objective-based amortization loss, corresponding to the
paper's Equation~(18).

\subsection*{Explanatory reconstruction}

\paragraph{Why is predicting a dual state natural?}
A first-year reader may reasonably ask: why predict dual variables instead of the
transport plan itself?  The answer is that the dual variables are a compact,
solver-native representation of the same teacher problem.  The chain is
\begin{equation}
\text{teacher problem}
\Longrightarrow
\text{dual state }(f,g)
\Longrightarrow
\text{recovered plan }\pi^*.
\label{eq:metaot-representation-chain}
\end{equation}
If the map from problem instance $(a,b,C)$ to optimal dual state is learnable, then
one can warm-start the teacher solve from a much better location than the usual
zero initialization.

\paragraph{Why predict only one dual half?}
Equation~\eqref{eq:metaot-dual-map} is the hidden simplification that makes the
paper practical.  Because $g$ can be reconstructed from $f$ by the teacher's own
first-order condition, predicting both halves is often unnecessary.  So the paper
reduces the learning burden from
\begin{equation}
(a,b,C)\mapsto (f,g)
\qquad\text{to}\qquad
(a,b,C)\mapsto f.
\label{eq:metaot-single-half-idea}
\end{equation}
This is not just an architectural convenience.  It means the model is learning a
teacher-consistent latent state rather than a free pair of vectors.

\paragraph{Why is the loss objective-based rather than supervised?}
The paper's Equation~\eqref{eq:metaot-amortization-loss} does not require ground
truth optimal duals.  Instead of regressing onto stored teacher solutions, it says:
make the predicted dual good according to the teacher's own objective.  This is why
the method can train without an explicit database of exact dual labels.

\paragraph{Why does fine-tuning preserve semantic conservatism?}
After prediction, one may still run a standard Sinkhorn correction.  So the
full deployment loop is
\begin{equation}
(a,b,C)
\xrightarrow{\text{predict }} \widehat f_\theta
\xrightarrow{\text{recover }\widehat g}
\xrightarrow{\text{corrective Sinkhorn}}
\widehat\pi_{\mathrm{corr}}.
\label{eq:metaot-corrected-pipeline}
\end{equation}
That last step is what preserves semantic conservatism: the final answer still
comes from the teacher dynamics, only faster.

\subsection*{Implementation manual: what Claude or a student would need}

\paragraph{Inputs.}
For the discrete route, the inputs are repeated OT problems of the form
\begin{equation}
(a,b,C,\varepsilon),
\label{eq:metaot-inputs-discrete}
\end{equation}
where $a$ and $b$ are discrete marginals and $C$ is the cost matrix.  For the
continuous route, the inputs are measure pairs $(\mu,\nu)$ together with whatever
representation the chosen teacher method uses.

\paragraph{Teacher objects.}
The implementer must decide which teacher is retained:
\begin{enumerate}[label=(\roman*)]
  \item discrete entropic OT with Sinkhorn,
  \item or continuous W2-style OT with a convex-potential teacher.
\end{enumerate}
This choice determines the latent state that is learned.

\paragraph{Parameterization.}
The discrete model predicts one dual vector,
\begin{equation}
\widehat f_\theta = M_\theta(a,b,C),
\label{eq:metaot-mlp}
\end{equation}
where $M_\theta$ is typically an MLP over the chosen problem representation.  The
continuous model predicts a dual potential or its parameter vector.

\paragraph{Training loop.}
A minimal discrete Meta OT loop is:
\begin{algorithm}[ht]
\caption{Meta OT training loop}
\begin{algorithmic}[1]
\Require repeated OT problem distribution $\mathcal D$ over $(a,b,C)$, regularization $\varepsilon$, model $M_\theta$
\For{each iteration}
  \State Sample a batch $(a,b,C)\sim\mathcal D$
  \State Predict $\widehat f_\theta=M_\theta(a,b,C)$
  \State Recover $\widehat g$ from \eqref{eq:metaot-dual-map}
  \State Evaluate the teacher objective $J(\widehat f_\theta;a,b,C)$
  \State Backpropagate and update $\theta$
\EndFor
\end{algorithmic}
\end{algorithm}

\paragraph{Deployment loop.}
At deployment time:
\begin{enumerate}[label=(\arabic*)]
  \item receive a new OT problem $(a,b,C)$,
  \item predict $\widehat f_\theta$,
  \item recover $\widehat g$,
  \item optionally build the approximate plan from \eqref{eq:metaot-primal-recovery},
  \item run corrective Sinkhorn iterations until the desired marginal error is met.
\end{enumerate}

\paragraph{Implementation hazards.}
The main hazards are:
\begin{enumerate}[label=(\roman*)]
  \item the representation of $(a,b,C)$ can dominate performance,
  \item one must decide whether to optimize a pure objective-based loss or include supervised regression to teacher duals,
  \item the correction loop must still be monitored by a teacher-side error criterion,
  \item the continuous route inherits the fragility of the underlying W2GN-style teacher.
\end{enumerate}

\paragraph{Implementation contract.}
The learned object is a solver-native latent state, not the final transport answer; the
optimized object is the predicted dual potential; the teacher solve is auxiliary but
non-negotiable; and the deployed artifact is the corrected teacher output after the
warm start has been refined.

\paragraph{Common wrong implementations.}
The easiest wrong turns are:
\begin{enumerate}[label=(\roman*)]
  \item training the predictor and skipping the corrective solve,
  \item regressing to arbitrary teacher labels without checking whether the objective-based contract is being preserved,
  \item evaluating predictor loss only and never checking the retained teacher residuals after correction.
\end{enumerate}
The first mistake changes semantics entirely: it turns a retained-teacher method into a
free predictor.

\paragraph{Minimum correct recipe.}
The shortest safe recipe is: fix the teacher, predict one teacher-native latent state,
recover the complementary dual information, and always finish by running the corrective
teacher solve before exporting the result.

\paragraph{First diagnostic.}
On a tiny batch of OT problems, compare warm-start plus correction against zero-init plus
correction under the same teacher-side residual criterion.  If the corrected outputs are
not at least teacher-consistent, the amortization route is not yet safe.

\subsection*{Code snapshot}

This is one of the strongest code-backed papers in the repo's availability ledger.
That ledger and the implementation-fit note both treat Meta OT as a high-value
starting point for the retained-teacher route.

Students should still read that status carefully:
\begin{enumerate}[label=(\roman*)]
  \item official code exists,
  \item it is highly relevant for amortized solver-state prediction,
  \item but it still requires adaptation to the BayesFilter teacher object and TensorFlow/TFP implementation conventions.
\end{enumerate}

\subsection*{BayesFilter bridge}

Meta OT remains the strongest retained-teacher paper for BayesFilter if the goal is
semantic conservatism.  The gap is not conceptual; it is in the representation and
scale of the predicted latent solver state.

\section{UNOT (Geuter et al., 2025)}

\subsection*{What problem is the paper solving?}

UNOT also remains close to a retained discrete entropic OT teacher, but with a
more operator-style architecture that aims to generalize across discretizations.
The paper is less about one fixed image size or one fixed discrete grid, and more
about learning an OT-related operator that works across varying resolutions.

\subsection*{Paper reproduction}

The paper starts from a dual OT formulation of the form
\begin{equation}
\sup_{\phi,\psi}
\left\{
\int \phi\,d\mu + \int \psi\,d\nu
\ : \ \phi(x)+\psi(y)\le c(x,y)
\right\},
\label{eq:unot-dual-abstract}
\end{equation}
and in the discrete entropic setting uses the Sinkhorn parameterization of the dual
via scaling variables $u,v$.  The standard Sinkhorn teacher loop is
\begin{equation}
u^{\ell+1} = a ./ (K v^{\ell}),
\qquad
v^{\ell+1} = b ./ (K^{\top}u^{\ell+1}),
\label{eq:unot-sinkhorn}
\end{equation}
followed by plan recovery
\begin{equation}
\pi = \diag(u)K\diag(v),
\label{eq:unot-plan-from-scales}
\end{equation}
which is the substance of the paper's Algorithm~1.

UNOT predicts a dual-derived latent quantity,
\begin{equation}
S_\theta(\mu,\nu)\approx g,
\qquad
v = \exp(g/\varepsilon),
\label{eq:unot-dual-prediction}
\end{equation}
From this latent quantity one reconstructs a plan, schematically through
\begin{equation}
\pi_{ij} \propto \exp\!\left(\frac{f_i+g_j-c_{ij}}{\varepsilon}\right).
\label{eq:unot-plan-recovery}
\end{equation}
The paper also introduces a synthetic data generator,
\begin{equation}
z\sim\rho_z \mapsto G_\vartheta(z)=(\mu,\nu),
\label{eq:unot-generator}
\end{equation}
with the actual implementation given in the paper's Equation~(6), and trains the
operator against Sinkhorn-generated targets through an $L_2$ loss, corresponding
to Equation~(7) and Algorithm~2.

\subsection*{Explanatory reconstruction}

\paragraph{What does UNOT learn that ordinary Sinkhorn does not?}
Ordinary Sinkhorn solves one problem at a time.  UNOT tries to learn the map
\begin{equation}
(\mu,\nu) \longmapsto \text{dual OT state},
\label{eq:unot-learned-map}
\end{equation}
so that solving the next OT problem becomes much cheaper.  The important point is
that the learned output is still a teacher-native object: a dual potential or a
scaling variable, not a free final transport map.

\paragraph{Why predict the log-potential rather than the plan directly?}
The plan itself changes with resolution and is constrained by marginals.  The dual
potential or its log-scaling form is often smoother and more learnable across
resolutions.  The paper therefore predicts
\begin{equation}
g = \varepsilon \log v,
\label{eq:unot-log-scaling}
\end{equation}
rather than the full plan matrix.  This is the same retained-teacher logic as Meta
OT, but adapted to a variable-resolution neural-operator setting.

\paragraph{Why does the synthetic generator matter?}
UNOT does not rely only on a fixed finite training set of OT problems.  The paper
constructs a generator of discrete measures across varying resolutions.  This is
why the method can claim resolution robustness: the training distribution itself is
built to vary the discretization.

A student implementing a first pass therefore has to make a concrete design choice
that the paper's abstract presentation can hide: what family of OT problems will be
generated during training, and what parts of that generator are fixed versus
learned?  In practice the generator is not a decorative data-loader.  It defines
the distribution of teacher problems on which the neural operator is expected to
generalize.

\paragraph{Why is the quotient-space issue important?}
The dual solution is not unique in a naive vector sense.  Equivalent dual states
can represent the same plan.  The paper uses this to motivate learning in a way
that respects the dual non-uniqueness rather than treating one arbitrary dual
representative as the only correct label.  This is one of the subtle points that
makes the operator-learning problem more stable than directly regressing onto a raw
plan matrix.

A first-year reader should translate that abstract statement into an operational
warning:
\begin{quote}
if two different dual vectors produce the same transport plan, then a naive
supervised loss can punish a correct prediction just because it picked a different
representative of the same equivalence class.
\end{quote}
So the quotient-space issue is not philosophical decoration.  It changes what a
sensible target construction and loss can look like.

\subsection*{Implementation manual: what Claude or a student would need}

\paragraph{Inputs.}
UNOT assumes discrete measures on
varying-resolution grids or meshes,
\begin{equation}
\mu = \sum_{i=1}^{m} a_i\delta_{x_i},
\qquad
\nu = \sum_{j=1}^{n} b_j\delta_{y_j},
\label{eq:unot-discrete-inputs}
\end{equation}
together with a cost function and entropic regularization parameter $\varepsilon$.

\paragraph{Teacher objects.}
The teacher is discrete entropic OT solved by Sinkhorn.  The core latent states are
its dual scalings or their logarithms.

\paragraph{Parameterization.}
The learned model is an operator
\begin{equation}
S_\theta : (\mu,\nu) \mapsto g,
\label{eq:unot-operator-param}
\end{equation}
typically realized as an FNO-like neural operator over the chosen measure
representation.  A separate synthetic generator $G_\vartheta$ creates training OT
instances.

\paragraph{Training loop.}
A minimal UNOT training loop is:
\begin{algorithm}[ht]
\caption{UNOT training loop}
\begin{algorithmic}[1]
\Require cost $c$, regularization $\varepsilon$, synthetic measure generator $G_\vartheta$, operator $S_\theta$, Sinkhorn target generator $\tau_k$
\For{each iteration}
  \State Sample latent noise $z\sim\rho_z$
  \State Generate a problem $(\mu,\nu)=G_\vartheta(z)$
  \State Predict $g_\theta=S_\theta(\mu,\nu)$
  \State Run the target generator $\tau_k$ to obtain improved teacher-side dual targets
  \State Minimize the $L_2$ discrepancy between predicted and target dual data
  \State Update both $\theta$ and, when applicable, the generator parameters $\vartheta$
\EndFor
\end{algorithmic}
\end{algorithm}

Here the two paper-specific objects that most need translation are $\tau_k$ and the
target duals.  Operationally, $\tau_k$ should be read as
\begin{quote}
a teacher-side target-refinement routine built from $k$ Sinkhorn updates or an
equivalent teacher-improvement map.
\end{quote}
So the training story is not simply ``predict one dual vector from data.''  It is
``predict a teacher-native latent state, then compare it against a Sinkhorn-improved
target built from the same OT instance.''  That retained-teacher interpretation is
what keeps UNOT close to a conservative BayesFilter lane.

Concretely, dual nonuniqueness means the loss should reward agreement at the level
of the represented transport plan or equivalence class, not only agreement with one
arbitrarily chosen raw dual vector.  A useful practical reading is:
\begin{enumerate}[label=(\roman*)]
  \item the teacher problem produces a family of equivalent dual descriptions,
  \item $\tau_k$ chooses or refines one teacher-side representative of that family,
  \item the learner should be judged by whether its prediction leads to the right corrected OT solution, not only by exact coordinate-wise equality to one raw label.
\end{enumerate}
This is the operational content behind the quotient-space warning.

\paragraph{Deployment loop.}
At deployment time:
\begin{enumerate}[label=(\arabic*)]
  \item encode the discrete source and target measures in the operator's expected format,
  \item predict the dual latent quantity $g_\theta$,
  \item convert to scaling variables by \eqref{eq:unot-log-scaling},
  \item either reconstruct the approximate plan directly or use the prediction as a Sinkhorn warm start,
  \item if high accuracy is needed, run corrective Sinkhorn iterations.
\end{enumerate}

\paragraph{Implementation hazards.}
The main difficulties are:
\begin{enumerate}[label=(\roman*)]
  \item the method is naturally grid or resolution-centric, so particle-cloud use requires a nontrivial representation bridge,
  \item the synthetic measure generator is part of the training story and must be implemented coherently,
  \item teacher-target generation still depends on repeated Sinkhorn solves during training,
  \item the predicted dual is not unique, so the loss must respect the teacher's equivalence structure.
\end{enumerate}

\paragraph{Implementation contract.}
The learned object is a dual-derived latent state; the optimized object is the operator
that predicts that latent state; the teacher-side target refinement remains auxiliary but
required; and the deployed artifact is either a corrected Sinkhorn warm start or the
transport object reconstructed from it after respecting the teacher equivalence class.

\paragraph{Common wrong implementations.}
The main wrong turns are:
\begin{enumerate}[label=(\roman*)]
  \item predicting one arbitrary raw dual vector and treating exact coordinate agreement as the real target,
  \item skipping the teacher-side refinement and treating the raw predictor output as the final answer,
  \item ignoring the representation mismatch between regular grids and unordered particle clouds.
\end{enumerate}
The first two change semantics rather than only efficiency.

\paragraph{Minimum correct recipe.}
The shortest safe recipe is: define the synthetic problem generator, predict the
teacher-native latent state, compare it against a teacher-refined target that respects
nonuniqueness, and deploy only through a corrected Sinkhorn-style route.

A practical first loss rule is therefore: do not supervise one arbitrary raw dual
vector.  Instead, choose a teacher-refined target representative, run the corrective
teacher updates needed to make that representative meaningful, and judge success by the
corrected OT solution and residual contract rather than by coordinate-wise agreement
alone.

\paragraph{First diagnostic.}
On a tiny family of OT problems, verify that warm start plus correction reproduces the
teacher residual contract and improves on naive initialization.  If that check fails, the
predictor may still be fitting an arbitrary dual label rather than a useful equivalence
class of teacher states.

\subsection*{Code snapshot}

This is also one of the strongest code-backed papers in the repo's ledger.  The
existing source-availability note records an official repository and treats UNOT as a
high-value retained-teacher-adjacent route.

The practical status is:
\begin{enumerate}[label=(\roman*)]
  \item official code exists,
  \item it is highly relevant for discrete entropic OT and solver-latent prediction,
  \item but it is not BayesFilter-ready without a bridge from grid or resolution-centric inputs to unordered particle clouds.
\end{enumerate}

\subsection*{BayesFilter bridge}

The paper is serious for BayesFilter because it keeps close contact with an
entropic teacher.  The hard part is the representation bridge from regular-grid or
resolution-centric operator inputs to unordered particle clouds.

\section{What this chapter established}

A first-year reader should now know that retained-teacher papers are not trying to
learn a final transport answer.  They are trying to learn a \emph{better way to
solve the same teacher problem}.  That is their main virtue for a cautious
BayesFilter design.


\chapter{Advanced directions and compressed shelf}

\section{How to read this chapter}

The remaining papers are scientifically useful, but they are not the best first
teaching path.  They belong here because they fill in the shelf and sharpen the
comparison, not because they should all be studied first.

However, the right fix is not to leave them compressed.  This chapter should now be
read with the same teaching-note ambition as the stronger earlier chapters: each
algorithm gets its own explanation of the learned object, objective, practical
optimization route, deployment object, BayesFilter relevance, and unresolved gap.

The main difference from the earlier chapters is not that the pedagogy should be
lighter.  It is that the frontier papers themselves sometimes support shorter
mathematical derivations or weaker implementation claims.  So the structure should
be uniform even when the depth of the paper's own executable detail is not.

As you read each Chapter 5 algorithm, keep four questions fixed:
\begin{enumerate}[label=(\roman*)]
  \item What object is learned or optimized?
  \item Why does the objective or formulation have that form?
  \item What is actually done in training or numerical optimization?
  \item What gets deployed, and why should BayesFilter care?
\end{enumerate}
\section{Static direct-map variants}

These papers all keep the learned object close to a direct map, but each changes a
different part of the burden: sparsity priors, Jacobian-bearing density formulas,
geometric optimization, fixed-point potential recovery, or manifold geometry.
Their value in this note is that they show how rich the direct-map lane becomes
once one asks for more than the original OT-ICNN story.

\subsection*{Chen (2025), displacement-sparse neural OT}

\paragraph{What problem is the paper solving?}
Chen keeps the direct-map OT setting but asks for a map whose displacement field is
structurally sparse.  The problem is therefore not to change the object class, but
to bias the learned map toward simpler or more localized motion.

\paragraph{Paper reproduction.}
The paper augments the direct-map objective by a sparsity regularizer on the
displacement,
\begin{equation}
J_{\mathrm{sparse}}(T_\theta)
=
J_{\mathrm{OT}}(T_\theta) + \lambda \mathcal R(T_\theta(u)-u).
\label{eq:chen-sparse-objective}
\end{equation}
So the deployed object is still a transport map $T_\theta$, but the optimization
criterion now prefers maps whose movement field has additional structure.

\paragraph{Explanatory reconstruction.}
The clean way to read the paper is:
\begin{align}
\text{ordinary direct transport map}
&\Longrightarrow
\text{add prior belief about how mass should move}
\notag\\
&\Longrightarrow
\text{penalize displacements that violate that structural prior}.
\label{eq:chen-logic}
\end{align}
Pedagogically, this is one of the simplest examples of how a paper can keep the
same learned object while changing the inductive bias.

\paragraph{Implementation manual: what Claude or a student would need.}
\paragraph{Inputs.}
One needs source and target samples exactly as in an ordinary direct-map OT setup.

\paragraph{Trainable object.}
The learned object is still the direct transport map $T_\theta$.

\paragraph{Objective.}
Optimize the OT fit term together with the displacement regularizer
\eqref{eq:chen-sparse-objective}.

\paragraph{Training loop.}
Each iteration samples minibatches, evaluates the transport-fit term, evaluates the
sparsity penalty on $T_\theta(u)-u$, combines them with weight $\lambda$, and
updates $\theta$.

\paragraph{Deployment loop.}
At deployment time, one simply applies $T_\theta$ to fresh source samples.  The
regularizer matters only because it changed the map learned during training.

\paragraph{Implementation hazards.}
The main hazard is over-reading the regularizer: if the true transport is not sparse
in the intended sense, the inductive bias can improve interpretability while hurting
fidelity.

\paragraph{BayesFilter bridge.}
This paper matters mainly as a design lesson.  It shows how one might inject prior
structure into a transport map, but it does not by itself solve BayesFilter's need
for repeated observation-conditioned posterior updates.

\paragraph{What the paper does not settle.}
It does not show that the chosen sparsity prior matches weighted-particle filtering
semantics or that sparsity is the right structural bias for BayesFilter clouds.

\subsection*{GradNetOT (Chaudhari 2025)}

\paragraph{What problem is the paper solving?}
GradNetOT stays in the direct-map family but makes the density-transform burden
explicit.  The question is no longer only where samples move, but also how the map's
Jacobian changes density.

\paragraph{Paper reproduction.}
The map is parameterized directly,
\begin{equation}
T_\theta : \R^d \to \R^d,
\label{eq:gradnet-map}
\end{equation}
with density relation
\begin{equation}
p(x)=q(T_\theta(x))\left|\det J_{T_\theta}(x)\right|,
\label{eq:gradnet-jacobian}
\end{equation}
and quadratic objective
\begin{equation}
\min_{T_\theta} \E_{x\sim p}\left[\frac12\|x-T_\theta(x)\|^2\right].
\label{eq:gradnet-ot}
\end{equation}
So the learned object is still a direct map, but the map must now carry explicit
Jacobian information.

\paragraph{Explanatory reconstruction.}
The right pedagogical staircase is:
\begin{align}
\text{direct map learning}
&\Longrightarrow
\text{need density transformation law}
\notag\\
&\Longrightarrow
\text{must control Jacobian information explicitly}.
\label{eq:gradnet-logic}
\end{align}
That is why the method feels heavier than a plain sample-pushforward map learner.

\paragraph{Implementation manual: what Claude or a student would need.}
\paragraph{Inputs.}
One needs samples together with a representation that makes the Jacobian-based terms
computable or approximable.

\paragraph{Trainable object.}
The learned object is still $T_\theta$.

\paragraph{Objective.}
The optimization combines transport cost with the density-transform contract coming
from \eqref{eq:gradnet-jacobian}.

\paragraph{Training loop.}
Each iteration evaluates the map, computes the relevant Jacobian term, forms the
training objective, and updates the network by backpropagation.

\paragraph{Deployment loop.}
Deployment still applies $T_\theta$ as a direct map, but the model is only credible
because training enforced the density-bearing structure.

\paragraph{Implementation hazards.}
Stable Jacobian evaluation is the main difficulty.  This is exactly why the paper is
more mathematically and numerically demanding than a simpler direct-map baseline.

\paragraph{BayesFilter bridge.}
The paper is a useful warning about how quickly direct-map transport becomes heavier
when exact density logic is made explicit.  That is informative for BayesFilter, but
it is not yet a natural first lane for repeated weighted-particle updates.

\paragraph{What the paper does not settle.}
It does not settle whether a Jacobian-heavy direct-map route is the right trade-off
for BayesFilter compared with more conditional or retained-teacher alternatives.

\subsection*{Dumont et al. (2026)}

\paragraph{What problem is the paper solving?}
Dumont keeps the learned object map-like, but changes the optimization story.  The
question becomes how to optimize transport maps geometrically rather than treating
them as only another unconstrained neural parameterization.

\paragraph{Paper reproduction.}
The paper optimizes
\begin{equation}
\argmin_{T_\theta\in K_{\mu}} D(\pushf{T_\theta}{\mu}\mid \nu),
\label{eq:dumont-divergence}
\end{equation}
with convex map
\begin{equation}
T_\theta = \nabla \phi_\theta,
\label{eq:dumont-convex-map}
\end{equation}
and gradient-flow dynamics
\begin{equation}
\partial_t T_t = -\operatorname{grad}\mathcal F(T_t).
\label{eq:dumont-gradient-flow}
\end{equation}
So the object is still a map, but the optimization route is geometric.

\paragraph{Explanatory reconstruction.}
The clean reading is:
\begin{align}
\text{transport map family}
&\Longrightarrow
\text{treat the family itself as a geometric space}
\notag\\
&\Longrightarrow
\text{optimize by a structured gradient flow on that space}.
\label{eq:dumont-logic}
\end{align}
This is why the paper feels different from both OT-ICNN and more standard neural map
fitting.

The phrase "geometric descent" can otherwise remain too vague.  The real shift is
that the update variable is not only a parameter vector sitting behind a neural map;
it is the map object itself viewed as living in a structured space.  So the descent
step is meant to respect that geometry rather than treat every admissible transport
map as just another unconstrained point in parameter space.

\paragraph{Implementation manual: what Claude or a student would need.}
\paragraph{Inputs.}
One needs source and target measures together with a choice of divergence
$D(\pushf{T_\theta}{\mu}\mid\nu)$.

\paragraph{Trainable object.}
The learned object is the convex-potential parameterization $\phi_\theta$, hence the
map $T_\theta=\nabla\phi_\theta$.

\paragraph{Objective.}
The target is not only to fit a map but to descend the divergence by the geometric
flow \eqref{eq:dumont-gradient-flow}.

\paragraph{Training or optimization loop.}
Each step evaluates the current pushforward mismatch, computes the geometric descent
update, and moves the map parameters along that structured direction.  In practice,
the student should think of this as repeatedly converting a current map into a
geometry-aware descent direction, then updating the potential representation so that
the next map is better with respect to the chosen divergence.

\paragraph{Deployment loop.}
Deployment applies the resulting map $T_\theta$ to fresh source samples.  So despite
the richer optimization story, the exported object is still a transport map rather
than a solver state or full path object.

\paragraph{Implementation hazards.}
The main burden is the optimization geometry itself: one must understand not just
how to parameterize the map, but how the structured descent is represented and
stabilized numerically.

\paragraph{BayesFilter bridge.}
The paper is scientifically valuable because it shows another way direct-map OT can
be regularized and optimized.  But it still reads as a pair-specific transport lane,
not yet as a repeated posterior-update block.

\paragraph{What the paper does not settle.}
It does not show that the geometric optimization burden pays off for BayesFilter's
repeated conditioned-update setting.

\subsection*{Park et al. (2026)}

\paragraph{What problem is the paper solving?}
Park focuses on recovering OT maps from a single potential by emphasizing the
$c$-transform structure.  The paper is therefore about making the solver side of the
potential story cleaner rather than changing the deployed object away from a map.

\paragraph{Paper reproduction.}
The paper emphasizes the $c$-transform,
\begin{equation}
f^c(y)=\inf_x\left\{\frac12\|x-y\|^2-f(x)\right\},
\label{eq:park-c-transform}
\end{equation}
and develops a fixed-point route for recovering the transport map from one
potential.

\paragraph{Explanatory reconstruction.}
The pedagogical lesson is:
\begin{align}
\text{single potential description}
&\Longrightarrow
\text{recover OT structure through the }c\text{-transform}
\notag\\
&\Longrightarrow
\text{shift burden from free map fitting to structured solver recovery}.
\label{eq:park-logic}
\end{align}
So the paper stays near the Brenier-style family while changing how the numerical
route is organized.

\paragraph{Implementation manual: what Claude or a student would need.}
\paragraph{Inputs.}
One needs a cost function and a potential parameterization that supports the
$c$-transform calculations.

\paragraph{Trainable or optimized object.}
The central object is the potential $f$, not an arbitrary direct map.

\paragraph{Objective.}
Optimization is driven by making the potential consistent with the fixed-point or
solver structure that recovers the map.

\paragraph{Training or optimization loop.}
Each step updates the potential representation and re-evaluates the fixed-point or
potential-recovery relation that defines the corresponding transport map.  So the
iterated variable is the potential, but the quantity one repeatedly checks is the map
recovered from that potential through the solver structure.

\paragraph{Deployment loop.}
Deployment evaluates the recovered map induced by the learned or optimized
potential.  The important practical message is that the paper does not deploy the raw
potential as the final user-facing object; it deploys the transport map obtained after
the fixed-point or potential-recovery step has been carried out.

\paragraph{Implementation hazards.}
The main hazard is that the potential-recovery logic can be elegant mathematically
while still being nontrivial to stabilize numerically.

\paragraph{BayesFilter bridge.}
The paper matters because it shows that direct-map OT can be made more solver-aware.
But the route is still not naturally phrased as a repeated observation-conditioned
update operator.

\paragraph{What the paper does not settle.}
It does not settle whether a cleaner fixed-point solver for a static OT potential is
the right abstraction for BayesFilter's actual transport needs.

\subsection*{Micheli et al. (2026)}

\paragraph{What problem is the paper solving?}
Micheli keeps the map-like object but moves the whole story to manifold geometry.
The problem is how to define and learn OT maps when Euclidean geometry is no longer
the right background.

\paragraph{Paper reproduction.}
The manifold cost is
\begin{equation}
c(x,y)=\frac12 d_{\mathcal M}(x,y)^2,
\label{eq:micheli-cost}
\end{equation}
and the deployed map is
\begin{equation}
T(x)=\operatorname{Exp}_x(-\nabla\phi(x)).
\label{eq:micheli-map}
\end{equation}
The deployed object is therefore still a map, but it lives inside manifold geometry
through the exponential map.

\paragraph{Explanatory reconstruction.}
The right staircase is:
\begin{align}
\text{OT map on Euclidean space}
&\Longrightarrow
\text{replace Euclidean geometry by manifold geometry}
\notag\\
&\Longrightarrow
\text{replace ordinary displacement by gradient plus exponential map}.
\label{eq:micheli-logic}
\end{align}
So the object remains familiar while the ambient geometry becomes much richer.

A slower way to read the formula is to compare it directly with the Euclidean story.
In flat space, a Brenier-style map says: take a gradient direction and move along it.
On a manifold, the gradient still tells us the infinitesimal direction of motion, but
that direction lives in the tangent space at the current point.  One cannot simply add
that tangent vector to the point itself.  The exponential map in
\eqref{eq:micheli-map} is the mechanism that turns the tangent-space direction back
into an actual point on the manifold.

A tiny concrete example helps.  If the state lived on a sphere, then a Euclidean
straight-line update would usually leave the sphere.  The manifold update instead says:
compute the tangent-space direction from the potential, then move along the sphere's
own geometry by the exponential map.  That is the smallest intuition behind the paper:
keep the transport step intrinsic to the geometry of the state space.

\paragraph{Implementation manual: what Claude or a student would need.}
\paragraph{Inputs.}
One needs source and target data on the manifold $\mathcal M$ together with the
manifold distance and differential operators.

\paragraph{Trainable object.}
The central trainable object is the potential $\phi$, whose gradient is interpreted
through the manifold geometry.

\paragraph{Objective.}
Optimization is driven by the manifold OT cost rather than the flat Euclidean one.

\paragraph{Training or optimization loop.}
Each iteration computes the geometry-aware objective, differentiates through the
manifold operators, and updates the potential representation.  In practice this means
that every optimization step must repeatedly move between three levels: the potential,
its tangent-space gradient, and the geometry-aware map obtained after applying the
exponential map.

\paragraph{Deployment loop.}
Deployment evaluates the map through the exponential-map formula
\eqref{eq:micheli-map}.  So the deployed object is still a map, but it is a map that
must be evaluated intrinsically on the manifold rather than by ordinary Euclidean
vector arithmetic.

\paragraph{Implementation hazards.}
The burden shifts to geometry: manifold gradients, exponential maps, and stable
numerical representations of the ambient space.  For BayesFilter, the right mental
question is very concrete: is the state space genuinely curved enough that ordinary
Euclidean transport would be misleading?  If not, the manifold burden may be
scientifically elegant but operationally unnecessary.

\paragraph{BayesFilter bridge.}
This is best read as a geometry frontier marker.  It matters if BayesFilter's state
space is intrinsically non-Euclidean, but otherwise it is a contrast paper rather
than the immediate next lane.

\paragraph{What the paper does not settle.}
It does not settle whether the additional manifold burden is justified for the
current BayesFilter use cases.
\section{Solver-side and unbalanced comparisons}

These papers move the emphasis away from learning one explicit direct map and
toward changing the solver object itself: faster retained solves, unbalanced plans,
or low-rank factorizations.  The key thing to watch is not just the formula, but
which optimization variable changes.

\subsection*{Kassraie et al. (2024)}

\paragraph{What problem is the paper solving?}
Kassraie asks how much can be gained by redesigning the entropic OT solve itself
without changing the semantic target.  The emphasis is therefore on solver quality,
not on replacing the OT object with a new learned map.

\paragraph{Paper reproduction.}
The retained teacher remains the entropic OT plan
\begin{equation}
\pi_\varepsilon
\in
\argmin_{\pi\in\Pi(a,b)}
\langle C,\pi\rangle + \varepsilon \KL(\pi\,\|\,a\otimes b).
\label{eq:kassraie-entropic}
\end{equation}
So the optimized object is still the regularized OT plan, together with whatever
solver state is used to reach it more efficiently.

\paragraph{Explanatory reconstruction.}
The clean reading is:
\begin{align}
\text{same semantic OT target}
&\Longrightarrow
\text{different numerical route to that target}
\notag\\
&\Longrightarrow
\text{performance gain comes from the solver, not from a new transport object}.
\label{eq:kassraie-logic}
\end{align}
This is exactly why the paper belongs in a solver-side section rather than in a map-
learning section.

A first-year student will understand the paper much better if they separate three
levels that are easy to collapse together:
\begin{enumerate}[label=(\roman*)]
  \item the semantic target, namely the same entropic OT solution as before,
  \item the iterated numerical state, such as dual variables, scalings, or another solver-side internal state,
  \item the final artifact kept after convergence, namely the solved plan or the transport object derived from it.
\end{enumerate}
The paper's claim lives almost entirely at level (ii): it changes how the retained
teacher is solved, not what the retained teacher means.

\paragraph{Implementation manual: what Claude or a student would need.}
\paragraph{Inputs.}
One needs the usual entropic OT ingredients $(a,b,C,\varepsilon)$.

\paragraph{Optimized object.}
The optimized object is still the regularized plan $\pi_\varepsilon$ or the
associated solver state.

\paragraph{Objective.}
The objective remains the same entropic OT objective \eqref{eq:kassraie-entropic}.

\paragraph{Optimization loop.}
The numerical loop repeatedly updates the solver state until the retained entropic OT
problem is solved to the desired tolerance.  The key point is that the loop is still
judged by teacher-side criteria: residual error, marginal mismatch, or whatever
numerical certificate the retained entropic solver uses.  So the accelerated method is
not finished when parameters stop moving; it is finished when the retained OT solve
has actually met its stopping rule.

A more explicit computational picture is helpful.  A first implementation would store
some solver-side state, run one accelerated update of that state, evaluate the same
teacher-side stopping diagnostic used by the baseline solver, and then either stop or
iterate again.  So the method still lives inside the same residual-based contract as
the teacher; it just tries to reach the acceptable residual regime faster.

\paragraph{Deployment loop.}
Deployment uses the resulting plan or barycentric projection exactly as in the
teacher method, because the semantic target never changed.  Operationally, the only
artifact that survives is the solved teacher object itself: the accelerated solver
state is valuable only because it helped reach that final retained OT solution more
quickly.

\paragraph{Implementation hazards.}
The main risk is over-claiming: a faster solver is not a new OT model, and it does
not by itself change the correctness contract of the retained teacher.

\paragraph{BayesFilter bridge.}
This is relevant whenever BayesFilter wants to keep the same OT teacher but reduce
solve cost.  It is therefore conceptually close to retained-teacher thinking.

\paragraph{What the paper does not settle.}
It does not by itself produce an observation-conditioned posterior map or an
amortized update rule.

\subsection*{Gazdieva et al. (2024)}

\paragraph{What problem is the paper solving?}
Gazdieva asks what happens when exact mass conservation is too restrictive.  The
problem is therefore not only transport, but transport with creation or destruction
of mass.

\paragraph{Paper reproduction.}
The unbalanced plan-side objective is
\begin{equation}
\inf_{\gamma}
\left\{
\int c(x,y)\,d\gamma
+\varepsilon\mathcal H(\gamma)
+\tau_1 D(\pi_1\gamma,\mu)
+\tau_2 D(\pi_2\gamma,\nu)
\right\},
\label{eq:gazdieva-ueot}
\end{equation}
and the parametric coupling factorization is
\begin{equation}
\gamma_\theta(x,y)=\mu_\omega(x)\gamma_\theta(y\mid x).
\label{eq:gazdieva-factorization}
\end{equation}
So the optimized or learned object is an unbalanced plan or sampler, not an exact
balanced OT coupling.

\paragraph{Explanatory reconstruction.}
The right staircase is:
\begin{align}
\text{balanced OT requires exact marginal match}
&\Longrightarrow
\text{relax the marginal constraints}
\notag\\
&\Longrightarrow
\text{optimize an unbalanced plan instead of a strictly balanced one}.
\label{eq:gazdieva-logic}
\end{align}
That is the conceptual move that changes the object class.

\paragraph{Implementation manual: what Claude or a student would need.}
\paragraph{Inputs.}
One needs source and target measures, a cost function, entropy regularization, and
penalties for marginal mismatch.

\paragraph{Trainable or optimized object.}
The core object is the coupling $\gamma$ or its parametric factorization.

\paragraph{Objective.}
Optimize the unbalanced OT objective \eqref{eq:gazdieva-ueot}, including the two
marginal-divergence penalties.

\paragraph{Training or optimization loop.}
Each iteration updates the coupling or factorization parameters to trade off cost,
entropy, and marginal mismatch.  Operationally, the loop is always balancing two
questions at once: are we moving mass cheaply, and are we allowing exactly the right
amount of mass creation or destruction for the application rather than simply washing
out the marginal constraints?

\paragraph{Deployment loop.}
Deployment uses the learned or optimized coupling for sampling or plan evaluation,
with the understanding that the result need not conserve mass exactly.  In a filtering
interpretation, this means the coupling is allowed to downweight or diffuse parts of
the source mass rather than insisting that every bit of predictive mass be matched one-
for-one to posterior mass.

\paragraph{Implementation hazards.}
The main hazard is semantic mismatch: once mass balance is relaxed, one must be very
clear whether that is scientifically appropriate for the target application.  For
BayesFilter, this is exactly why the paper remains a non-default lane: unbalanced OT
may be useful in some modeling stories, but it changes the meaning of the update more
radically than the balanced retained-teacher or conditional-map routes do.

\paragraph{BayesFilter bridge.}
This is relevant if the filtering story genuinely wants unbalanced transport.  It is
not the default BayesFilter lane, but it is an important contrast when strict mass
preservation is not the right abstraction.

\paragraph{What the paper does not settle.}
It does not show that unbalanced transport is the right default semantics for the
current BayesFilter particle-update problem.

\subsection*{Scetbon et al. (2023)}

\paragraph{What problem is the paper solving?}
Scetbon asks how to keep the plan-side object but compress its numerical
representation.  The problem is computational: how to avoid carrying a dense
coupling matrix when a low-rank structure may suffice.

\paragraph{Paper reproduction.}
The low-rank unbalanced OT formulation is
\begin{equation}
\mathrm{ULOT}_r(a,b)
=
\min_{\pi:\,\operatorname{rank}_+(\pi)\le r}
\langle C,\pi\rangle
+\tau_1\KL(\pi\mathbf 1\mid a)
+\tau_2\KL(\pi^{\top}\mathbf 1\mid b),
\label{eq:scetbon-ulot}
\end{equation}
with low-rank factorization
\begin{equation}
\pi = Q\diag(1/g)R^{\top}.
\label{eq:scetbon-low-rank-factor}
\end{equation}
So the object class stays plan-side, but its numerical representation changes
fundamentally.

\paragraph{Explanatory reconstruction.}
The pedagogical lesson is:
\begin{align}
\text{dense plan matrix}
&\Longrightarrow
\text{replace by low-rank factors}
\notag\\
&\Longrightarrow
\text{computational savings come from representation, not from a different semantic target}.
\label{eq:scetbon-logic}
\end{align}
That is why the paper belongs beside solver-side methods.

A student should make the factorization in \eqref{eq:scetbon-low-rank-factor}
concrete.  In a dense plan, every source location can in principle talk to every
target location through one large matrix $\pi$.  In the low-rank picture,
\begin{equation}
\pi = Q\diag(1/g)R^{\top},
\end{equation}
that communication is forced to pass through a much smaller latent bottleneck.  One
can think of the columns of $Q$ and $R$ as describing a small collection of latent
transport channels, while $g$ rescales how strongly those channels are used.  So the
plan is no longer stored entry-by-entry; it is reconstructed from a compressed set of
factors.

This also explains the main benefit and the main risk.  The benefit is obvious:
storing and updating $Q$, $g$, and $R$ can be much cheaper than storing a dense
matrix.  The risk is equally important: if the chosen rank is too small, the factor
family may simply be unable to represent the transport geometry the problem really
needs.  Then the method fails not because OT is wrong, but because the compression is
too aggressive.

\paragraph{Implementation manual: what Claude or a student would need.}
\paragraph{Inputs.}
One needs the usual OT ingredients plus a chosen rank budget $r$.

\paragraph{Optimized object.}
The optimized object is the low-rank factorization $(Q,g,R)$ representing the plan.

\paragraph{Objective.}
Optimize the low-rank unbalanced OT objective \eqref{eq:scetbon-ulot}.

\paragraph{Optimization loop.}
Each step updates the factor matrices and scaling variables rather than a full dense
plan.  Operationally, the loop repeatedly improves the latent factors, reconstructs
the implied coupling, checks the transport objective and marginal-penalty terms, and
continues until the compressed representation is good enough for the target error
budget.

\paragraph{Deployment loop.}
Deployment reconstructs or samples from the plan through the low-rank
representation.  If BayesFilter wanted transported particles rather than only a plan,
one would still need the usual final step of turning the recovered coupling into a
barycentric or sampling-based particle update.  So the low-rank compression changes
how the coupling is represented, not the basic logic by which a coupling is turned
into moved particles.

\paragraph{Implementation hazards.}
The key hazard is choosing a rank that is computationally cheap without destroying
the transport fidelity one actually needs.  A practical symptom of rank failure would
be a plan that is fast and memory-light, but that cannot reproduce the necessary
transport structure even when the numerical optimization appears to converge.

\paragraph{BayesFilter bridge.}
This is important whenever BayesFilter wants to keep plan-side transport but make it
more scalable.  It is therefore a computational bridge, not yet a new posterior-map
family.

\paragraph{What the paper does not settle.}
It does not show that low-rank plan structure is the right abstraction for the
weighted-cloud geometries BayesFilter will actually encounter.
\section{Dynamic paths and measure operators}

This section moves from endpoint maps or plans to richer objects: whole paths of
measures, measure-to-measure operators, or local families of transport rules.  That
is why it is the most ambitious frontier section in the chapter.

The easiest way to read it is in two passes.  First, use GeONet as the anchor for
what it means to learn a path-valued or operator-valued OT object.  Then read the
later papers as nearby variants of that larger object class: one emphasizes motion
fields, one emphasizes measure-to-measure operators, and one emphasizes local rather
than global transport rules.

\subsection*{GeONet (2024)}

\textbf{What problem is the paper solving?}
GeONet is not trying to learn a final transport map.  It is trying to learn the
entire Wasserstein geodesic between two endpoint measures.  So the learned object
is a \emph{path of distributions}, not a single map.  If students do not start from
that sentence, the rest of the section becomes unreadable.

A useful mental model is a movie.  At time $t=0$ we have one density $\mu_0$; at
time $t=1$ we have another density $\mu_1$.  GeONet is trying to learn the whole
movie
\begin{equation}
(\mu_t)_{t\in[0,1]},
\label{eq:geonet-movie}
\end{equation}
not merely the last frame.  That is why the paper is harder than ordinary map
learning and why a PDE appears almost immediately.

A tiny 1D picture is helpful.  Imagine that $\mu_0$ is concentrated near the left
side of a line and $\mu_1$ near the right side.  Then the learned object is not one
narrow arrow from left to right.  It is the whole family of intermediate snapshots
showing how mass moves across the line over time.  In that picture,
\begin{itemize}
  \item $\rho_t$ tells us what the density snapshot looks like at time $t$,
  \item $v_t$ tells us how mass is moving at that time,
  \item and $u$ is the potential whose gradient generates that motion.
\end{itemize}
This is the smallest concrete interpretation of the ``movie'' language.

\textbf{Population dynamic OT problem.}
The path is not arbitrary.  It is chosen by the dynamic OT or Benamou--Brenier
principle:
\begin{equation}
\frac12\W_2^2(\mu_0,\mu_1)
=
\inf_{(\rho_t,v_t)}
\int_0^1 \int \frac12\|v_t(x)\|^2\rho_t(x)\,dx\,dt
\quad\text{subject to}\quad
\partial_t\rho_t + \nabla\cdot(\rho_t v_t)=0.
\label{eq:geonet-benamou-brenier}
\end{equation}
This is the real mathematical heart of the paper.  It says: among all paths of
densities connecting the endpoints, choose the one with smallest total kinetic
energy.

\textbf{Why does a PDE appear at all?}
If a density changes continuously in time, mass cannot disappear or appear
arbitrarily.  It must flow.  That is exactly the continuity equation,
\begin{equation}
\partial_t\rho_t + \nabla\cdot(\rho_t v_t)=0.
\label{eq:geonet-continuity-base}
\end{equation}
In words: density changes only because mass moves through space.  This is the first
PDE students must internalize before anything about KKT systems is mentioned.

\textbf{How the KKT system arises.}
To optimize \eqref{eq:geonet-benamou-brenier}, one introduces a Lagrange
multiplier $u(x,t)$ for the continuity equation.  The Lagrangian is
\begin{equation}
\mathcal L(\rho,v,u)
=
\int_0^1\int
\left[
\frac12\|v\|^2\rho + (\partial_t\rho + \nabla\cdot(\rho v))u
\right]dx\,dt,
\label{eq:geonet-lagrangian}
\end{equation}
which is the paper's Equation~(5).  Stationarity with respect to the velocity field
shows that the optimal velocity must satisfy
\begin{equation}
v_t^*(x)=\nabla u^*(x,t),
\label{eq:geonet-optimal-velocity}
\end{equation}
and stationarity with respect to the density yields the Hamilton--Jacobi equation
\begin{equation}
\partial_t u + \frac12\|\nabla u\|^2 = 0.
\label{eq:geonet-hj}
\end{equation}
Substituting the optimal velocity back into the continuity equation gives
\begin{equation}
\partial_t\rho_t + \nabla\cdot(\rho_t\nabla u)=0.
\label{eq:geonet-continuity}
\end{equation}
Together these form the paper's KKT system
\begin{equation}
\begin{cases}
\partial_t\rho_t + \nabla\cdot(\rho_t\nabla u)=0,\\
\partial_t u + \dfrac12\|\nabla u\|^2=0,\\
\rho(\cdot,0)=\mu_0,\qquad \rho(\cdot,1)=\mu_1.
\end{cases}
\label{eq:geonet-kkt}
\end{equation}
This is the conceptual staircase that the original paper expects the reader to
reconstruct.  The note now makes it explicit:
\begin{align}
\text{dynamic OT objective}
&\Longrightarrow
\text{continuity constraint}
\notag\\
&\Longrightarrow
\text{Lagrange multiplier}
\notag\\
&\Longrightarrow
\text{velocity and HJ equations}.
\label{eq:geonet-logic-chain}
\end{align}

\textbf{What exactly is GeONet learning?}
The paper's learned object is the nonlinear operator
\begin{equation}
\Gamma^\dagger : \calP_2(\Omega)\times\calP_2(\Omega) \to AC(\calP_2(\Omega)),
\qquad
(\mu_0,\mu_1)\mapsto (\mu_t)_{t\in[0,1]},
\label{eq:geonet-operator}
\end{equation}
which matches the paper's Equations~(12)--(13).  This is why the network cannot be
just an endpoint map regressor.  It must output a whole path or enough structure to
evaluate that path at any desired time.

\textbf{Parameterization.}
The paper introduces coupled DeepONet-style parameterizations for the density path
and the dual potential.  In the note's notation we can write the two operator
outputs as
\begin{equation}
C_\phi(\mu_0,\mu_1)(x,t)
\approx \rho_t(x),
\qquad
H_\psi(\mu_0,\mu_1)(x,t)
\approx u(x,t),
\label{eq:geonet-two-outputs}
\end{equation}
corresponding to the paper's Equations~(14)--(15).  The first network tries to
represent the path of densities; the second tries to represent the potential whose
gradient gives the optimal velocity.

\textbf{Why the losses look the way they do.}
Because the target object is the KKT system, the loss is built from residuals of
that system.  The continuity residual loss is
\begin{equation}
L_{\mathrm{cty}}
=
\alpha_1\E\left[\left|\partial_t C_\phi + \nabla\cdot(C_\phi\nabla H_\psi)\right|^2\right],
\label{eq:geonet-lcty}
\end{equation}
and the Hamilton--Jacobi residual loss is
\begin{equation}
L_{\mathrm{HJ}}
=
\alpha_2\E\left[\left|\partial_t H_\psi + \frac12\|\nabla H_\psi\|^2\right|^2\right].
\label{eq:geonet-lhj}
\end{equation}
Boundary losses enforce the endpoint densities:
\begin{equation}
L_{\mathrm{BC}}
=
\beta_0\E\bigl[\|C_\phi(\cdot,0)-\mu_0\|_{L^2}^2\bigr]
+
\beta_1\E\bigl[\|C_\phi(\cdot,1)-\mu_1\|_{L^2}^2\bigr],
\label{eq:geonet-lbc}
\end{equation}
which are the paper's Equations~(17)--(19) in compressed notation.  The total
training objective is then
\begin{equation}
\min_{\phi,\psi} L_{\mathrm{cty}}+L_{\mathrm{HJ}}+L_{\mathrm{BC}}.
\label{eq:geonet-total-loss}
\end{equation}
Pedagogically, this is the right way to read the paper: each term corresponds to a
specific failure mode of the learned path.

\textbf{Implementation manual: what Claude or a student would need.}

\paragraph{Inputs.}
The training data are pairs of endpoint measures,
\begin{equation}
(\mu_0^{(i)},\mu_1^{(i)}),
\qquad i=1,\ldots,n,
\label{eq:geonet-endpoint-dataset}
\end{equation}
represented on some branch grid or mesh.  One must also sample collocation points
$(x,t)$ in space-time to evaluate PDE residuals.

\paragraph{Trainable objects.}
The code needs two coupled outputs:
\begin{enumerate}[label=(\roman*)]
  \item a density-path network $C_\phi$,
  \item a potential network $H_\psi$.
\end{enumerate}
Both depend on the endpoint pair $(\mu_0,\mu_1)$ and the query point $(x,t)$.

\paragraph{Training loop.}
A minimal implementation loop is:
\begin{algorithm}[ht]
\caption{GeONet training loop}
\begin{algorithmic}[1]
\Require endpoint-measure dataset, branch grid, collocation sampler, networks $C_\phi,H_\psi$
\For{each iteration}
  \State Sample endpoint pairs $(\mu_0,\mu_1)$
  \State Sample collocation points $(x,t)$ in space-time
  \State Evaluate $C_\phi(\mu_0,\mu_1)(x,t)$ and $H_\psi(\mu_0,\mu_1)(x,t)$
  \State Compute continuity residual \eqref{eq:geonet-lcty}
  \State Compute HJ residual \eqref{eq:geonet-lhj}
  \State Compute endpoint or boundary residual \eqref{eq:geonet-lbc}
  \State Backpropagate the total loss \eqref{eq:geonet-total-loss}
  \State Update $\phi,\psi$
\EndFor
\end{algorithmic}
\end{algorithm}
This is the operational meaning of the paper's Algorithm~1.

\paragraph{Deployment loop.}
At deployment time:
\begin{enumerate}[label=(\arabic*)]
  \item provide two endpoint measures $(\mu_0,\mu_1)$,
  \item evaluate $C_\phi$ and optionally $H_\psi$ on a desired mesh and times $t$,
  \item interpret the output $C_\phi$ as the learned geodesic path,
  \item if needed, reconstruct a velocity field from $\nabla H_\psi$.
\end{enumerate}
So the deployed object is a path evaluator, not a one-shot endpoint transport map.

\paragraph{Implementation hazards.}
The main dangers are:
\begin{enumerate}[label=(\roman*)]
  \item the method learns densities or density-like fields, not just particles, so a mesh or discretization choice is unavoidable,
  \item PDE residuals require stable automatic differentiation through coupled operator outputs,
  \item endpoint conditions must be enforced strongly enough or the path can look plausible while missing the actual endpoints,
  \item the learned object is much richer than what BayesFilter currently needs, so successful training does not imply an easy insertion point into the existing filtering code.
\end{enumerate}

\paragraph{Implementation contract.}
The learned objects are the path network $C_\phi$ and the potential network $H_\psi$;
the optimized object is the coupled PDE-residual training problem; the KKT system is an
auxiliary correctness certificate; and the deployed object is the evaluated path (and
optionally velocity field), not a one-shot endpoint map.

\paragraph{Common wrong implementations.}
The most likely wrong turns are:
\begin{enumerate}[label=(\roman*)]
  \item training only for endpoint fit and ignoring continuity or HJ residuals,
  \item exporting one endpoint sample as if the paper learned only a terminal map,
  \item treating the dual potential $u$ as the deployed object instead of the path or the velocity induced by it.
\end{enumerate}
The first two change the algorithm's contract, not just its efficiency.

\paragraph{Minimum correct recipe.}
The shortest safe recipe is: represent endpoint measures on a finite mesh, train the two
coupled operator outputs against the continuity, HJ, and boundary losses, and deploy by
evaluating the path (and optional velocity) over the desired time grid.

\paragraph{First diagnostic.}
Before any broad experiment, test whether the learned path hits the endpoints and
reduces the continuity and HJ residuals on held-out collocation points.  A path that
looks visually smooth but misses those two checks is not yet implementing the paper's
contract.  The safest tiny first build is a 1D Gaussian-to-Gaussian path on a fixed
mesh, because then one can see immediately whether the learned movie, endpoint fit,
and residual logic are all behaving coherently.

\textbf{What the paper leaves implicit.}
The original paper moves very quickly from the Benamou--Brenier formulation to the
PDE residual training story.  The missing bridge is exactly the one many students
need:
\begin{align}
\text{OT path problem}
&\Longrightarrow
\text{mass conservation constraint}
\notag\\
&\Longrightarrow
\text{Lagrange multiplier}
\notag\\
&\Longrightarrow
\text{coupled PDE residuals as training losses}.
\label{eq:geonet-teaching-bridge}
\end{align}
Once that ladder is made explicit, the role of the KKT system and of the operator
network becomes much easier to understand.

\textbf{BayesFilter bridge.}
GeONet matters if BayesFilter eventually wants to learn a whole transport path or an
operator-valued continuation of OT rather than only one posterior update.  That is
scientifically exciting, but it is a larger object than BayesFilter immediately
needs.

\textbf{What the paper does not settle.}
GeONet does not by itself tell us how to insert a path-valued density operator into
the current weighted-particle filtering loop without an additional representation
bridge.

\subsection*{Chen and Liu (2025)}

\paragraph{What problem is the paper solving?}
Chen and Liu also move to dynamic transport, but from a more sample-based or
variational solver viewpoint.  The question is how to learn the motion field that
carries mass, rather than only the final map.

\paragraph{Paper reproduction.}
The dynamic variational problem is
\begin{equation}
\inf_{(\rho_t,v_t)} \int_0^1 \int L(v_t(x))\rho_t(x)\,dx\,dt
\quad\text{subject to}\quad
\partial_t\rho_t + \nabla\cdot(\rho_t v_t)=0.
\label{eq:chenliu-dynamic}
\end{equation}
So the central object is again a path and velocity field, not only a terminal map.

\paragraph{Explanatory reconstruction.}
The simplest reading is:
\begin{align}
\text{endpoint transport problem}
&\Longrightarrow
\text{represent the whole motion of mass}
\notag\\
&\Longrightarrow
\text{learn or optimize the velocity field that produces that motion}.
\label{eq:chenliu-logic}
\end{align}
This makes the paper a close frontier cousin of GeONet.

A student can get less lost by separating three objects that are easy to blur
together:
\begin{enumerate}[label=(\roman*)]
  \item the terminal pushforward or endpoint agreement,
  \item the full path of intermediate densities $\rho_t$,
  \item and the velocity field $v_t$ that actually drives that path.
\end{enumerate}
The paper is not only trying to get the endpoints right.  It is trying to learn or
optimize the motion law that takes us from the initial state to the final one.

In practice that means time has to be discretized.  A first implementation would not
store an abstract continuum of times; it would choose a time grid
$0=t_0<t_1<\cdots<t_K=1$, represent the path or velocity on that grid, and enforce
the continuity logic numerically step by step.  This is the most concrete bridge from
the variational formula to code.

\paragraph{Implementation manual: what Claude or a student would need.}
\paragraph{Inputs.}
One needs endpoint measures or sample representations together with a way to
represent the evolving density and motion field.

\paragraph{Trainable or optimized object.}
The central object is the path-density/velocity pair $(\rho_t,v_t)$ or its neural
surrogate.

\paragraph{Objective.}
Optimization is driven by the dynamic action plus the continuity constraint in
\eqref{eq:chenliu-dynamic}.

\paragraph{Training or optimization loop.}
Each step updates the path or velocity representation so that the dynamic objective
falls while the continuity constraint is respected numerically.  In a time-discretized
implementation, that means repeatedly updating the path or velocity on the chosen time
grid, evaluating the discrete action, and checking whether the induced transport over
successive time slices still respects the continuity logic.

\paragraph{Deployment loop.}
Deployment evaluates the learned motion or the resulting path, not merely one
terminal transport map.  For BayesFilter, the practical question would be whether the
filter consumes the full motion law, a sequence of intermediate transport states, or
only the final endpoint map obtained after integrating the learned velocity field.

\paragraph{Implementation hazards.}
The main hazards are stable dynamic optimization and a usable finite representation
of the path object.  A second hazard is interface mismatch: a path-based method may be
scientifically attractive, but a particle filter may only want a terminal update map,
so the student must decide what part of the dynamic object is actually exported.

\paragraph{BayesFilter bridge.}
This is important as a path-based frontier lane.  It is not yet the clearest first
implementation route for repeated observation-conditioned filtering updates.

\paragraph{What the paper does not settle.}
It does not settle how a learned path object should be integrated into BayesFilter's
current particle-update semantics.

\subsection*{Vandergrift et al. (2026)}

\paragraph{What problem is the paper solving?}
Vandergrift asks for a measure-to-measure operator that generalizes across many
transport instances.  The key shift is from learning one map for one pair to
learning an operator acting on whole measures.

\paragraph{Paper reproduction.}
The learned operator is
\begin{equation}
\mathcal G_\theta : \calP(\R^d)\to\calP(\R^d),
\qquad
\nu\approx \mathcal G_\theta(\mu),
\label{eq:vandergrift-operator}
\end{equation}
with static training objective
\begin{equation}
\min_\theta \sum_{i=1}^{m}\mathcal L\bigl(\mathcal G_\theta(\mu_i),\nu_i\bigr).
\label{eq:vandergrift-loss}
\end{equation}
The dynamic variant instead learns a measure-dependent velocity field.

\paragraph{Explanatory reconstruction.}
The pedagogical staircase is:
\begin{align}
\text{one transport map per pair}
&\Longrightarrow
\text{one operator that acts on many source measures}
\notag\\
&\Longrightarrow
\text{generalization target shifts from point samples to whole measures}.
\label{eq:vandergrift-logic}
\end{align}
That is why this is more ambitious than the ordinary direct-map story.

A student can make this concrete by imagining two different encodings of a source
measure: an unordered particle cloud and a fixed grid histogram.  The operator view
says that neither specific encoding is the mathematical object of interest.  The real
object is the measure-to-measure rule that should survive a reasonable change of
representation.  This is the practical meaning of ``generalize across measures'': not
memorize one map for one cloud, but learn a transformation law that accepts many
measure instances from the same family.

\paragraph{Implementation manual: what Claude or a student would need.}
\paragraph{Inputs.}
One needs a dataset of source and target measures together with a loss
$\mathcal L$ defined on measures or their representations.

\paragraph{Trainable object.}
The trainable object is the operator $\mathcal G_\theta$ or, in the dynamic case, a
measure-dependent velocity field.

\paragraph{Objective.}
Optimize the operator so that $\mathcal G_\theta(\mu_i)$ matches $\nu_i$ across a
family of measure pairs.

\paragraph{Training loop.}
Each iteration samples measure pairs, evaluates the operator, computes the measure-
level loss, and updates the operator parameters.  The hidden practical burden is that
the measures must first be encoded in some finite representation: particles, grids,
features, kernels, or another summary that the operator can actually process.

This is where the paper becomes more algorithmic than a slogan like ``learn an
operator on measures'' may suggest.  A real implementation must define three maps,
not one: an encoder from a concrete measure representation into the operator's input
space, the operator itself, and a decoder or output interpretation that turns the
operator result back into something usable.  Without those three pieces, the phrase
``measure-to-measure operator'' remains mathematically attractive but computationally
underspecified.

\paragraph{Deployment loop.}
Deployment applies $\mathcal G_\theta$ to a new source measure, rather than applying
one pair-specific map to one cloud.  So the deployed object is not "the map for this
problem" but a callable measure-to-measure transformation rule.  In BayesFilter
terms, this would look like an update block that accepts an entire predictive cloud or
measure representation and returns a transformed posterior-side representation.

\paragraph{Implementation hazards.}
The main burden is representation: measures must be encoded in a way that supports
operator learning without destroying the geometry one cares about.  A second hazard is
decoding: once the operator outputs a transformed measure representation, the student
must still decide how BayesFilter consumes it as particles, samples, or another usable
state object.

\paragraph{BayesFilter bridge.}
This is a natural frontier paper if BayesFilter eventually wants a true
measure-to-measure update operator.  That is broader than the current immediate
transport block, but still highly relevant as a long-run direction.

\paragraph{What the paper does not settle.}
It does not settle whether a global measure operator is preferable to a more local,
conditional, or retained-teacher BayesFilter route.

\subsection*{Girshfeld and Chen (2025)}

\paragraph{What problem is the paper solving?}
Girshfeld and Chen ask what happens if one does not trust one global transport rule.
The idea is to use locality in measure space so that nearby source measures select
or blend nearby transport rules.

\paragraph{Paper reproduction.}
The local transport family is
\begin{equation}
T_\theta[\mu] : \calP(\R^d)\to\calP(\R^d),
\label{eq:girshfeld-local-family}
\end{equation}
with local smoother
\begin{equation}
\widehat T_\theta[\mu_0]
=
\sum_{i=1}^{m} K_h\bigl(\W_2(\mu_0,\mu_i)\bigr)T_{\theta_i}.
\label{eq:girshfeld-local-smoother}
\end{equation}
So the learned object is not one universal operator, but a local family of transport
rules blended according to Wasserstein proximity.

\paragraph{Explanatory reconstruction.}
The clean reading is:
\begin{align}
\text{do not trust one global rule}
&\Longrightarrow
\text{store local rules in measure space}
\notag\\
&\Longrightarrow
\text{blend them according to Wasserstein locality}.
\label{eq:girshfeld-logic}
\end{align}
Pedagogically, this is a Wasserstein analogue of local regression.

A concrete picture helps here.  Suppose BayesFilter encounters one predictive cloud
that looks nearly Gaussian, another that is clearly multimodal, and a third that lies
between them.  The local-rule philosophy says that the third cloud should not be
forced to use the exact same transport logic as the first two.  Instead, it should
borrow mostly from nearby examples in measure space.  So locality is the paper's way
of weakening the assumption that one global operator is equally trustworthy
everywhere.

\paragraph{Implementation manual: what Claude or a student would need.}
\paragraph{Inputs.}
One needs a family of source measures, corresponding local transport rules, and a
kernel or neighborhood rule defined in Wasserstein space.

\paragraph{Trainable or optimized object.}
The central object is the family $T_\theta[\mu]$ together with the local smoothing
rule.

\paragraph{Objective.}
The optimization target is to learn useful local rules whose blend generalizes well
to nearby source measures.

\paragraph{Training loop.}
Each iteration updates the local transport family and the weighting or smoothing
scheme so that nearby measures receive compatible local transport behavior.

Operationally, that means the algorithm must repeatedly answer two questions at
once: what are the candidate local rules, and which training measures should count as
"nearby" enough to influence one another?  So the learned object is never only the
transport family by itself; it is the transport family together with the locality
mechanism that chooses how to blend it.

\paragraph{Deployment loop.}
Deployment takes a new source measure $\mu_0$, finds nearby reference measures in
Wasserstein space, and blends the corresponding local transport rules through
\eqref{eq:girshfeld-local-smoother}.

\paragraph{Implementation hazards.}
The main burdens are locality design, efficient Wasserstein-neighborhood search, and
ensuring the blended local rules remain numerically stable.

\paragraph{BayesFilter bridge.}
This is relevant if BayesFilter wants local rather than universal transport logic:
a different update rule for different regions of measure space rather than one
single global mechanism.

\paragraph{What the paper does not settle.}
It does not settle whether such locality is better than a single conditional update
operator or a retained-teacher route for the current BayesFilter setting.

\chapter{What should BayesFilter do next?}

\section{The main scientific lesson}

The papers on the shelf are not all answering the same question.  Once that is
understood, the main ranking becomes clearer.

\begin{enumerate}[label=(\arabic*)]
  \item If the question is \emph{conditional posterior transport}, the core papers
  are Wang, Al-Jarrah, and Hosseini.
  \item If the question is \emph{semantic conservatism with a retained teacher},
  the core papers are Meta OT and UNOT.
  \item If the question is \emph{advanced re-architecture}, the dynamic and
  measure-operator papers become relevant, but they should not be confused with the
  immediate BayesFilter transport problem.
\end{enumerate}

\section{The top two experimental lanes}

\paragraph{Lane A: conditional posterior transport.}
This is the scientifically most central lane.  Use Wang as the leading practical
paper, Al-Jarrah as the filtering-native contrast, and Hosseini as the strongest
theorem-level conditional OT anchor.

\paragraph{Lane B: retained-teacher acceleration.}
This is the more conservative lane.  Use Meta OT as the leading paper, UNOT as the
main alternative, and keep solver-side baselines such as Kassraie or Scetbon in
view.

\section{What should LEDH-PFPF-OT try first?}

The right recommendation for LEDH-PFPF-OT should start from the transport object it
already uses, not from the most fashionable paper on the shelf.  In the current
BayesFilter direction, LEDH-PFPF-OT is not an abstract Wasserstein comparison
problem.  It is a finite entropic-OT / Sinkhorn-style transport correction layer with
barycentric equal-weight output, implemented under the repository's declared default
production direction:
\begin{equation}
\text{TensorFlow/TFP},\qquad \text{GPU execution},\qquad \text{float32 with TF32 enabled},
\label{eq:ledhpfpfot-default-direction}
\end{equation}
with streaming or chunked execution where needed.  This owner directive appears in
`AGENTS.md` and `CLAUDE.md`.  It is an engineering target, not a proof of posterior
correctness or HMC readiness, but it fixes the implementation question we are trying
to answer.

The most important mathematical point is that LEDH-PFPF-OT already has a declared
teacher object: a finite entropic OT correction computed by Sinkhorn-style scaling and
exported through a barycentric equal-weight cloud.  So the ranking criterion should be:
\begin{quote}
prefer the candidate that preserves the current transport object and its executable
invariants before the candidate that replaces that object by a looser surrogate.
\end{quote}
The relevant invariants are:
\begin{enumerate}[label=(\roman*)]
  \item row and column marginal residuals for the finite coupling,
  \item the barycentric output semantics of the equal-weight cloud,
  \item the exact executed scalar/gradient path of the implementation,
  \item and compatibility with the GPU / JIT / TF32 execution target.
\end{enumerate}
A route that preserves these invariants is easier to justify as a next step than a
route that changes them, even if both are scientifically interesting.

Under that criterion, the ranking is:
\begin{enumerate}[label=(\arabic*)]
  \item exact-semantics entropic OT reference lane, realized by streaming/online/GPU Sinkhorn where possible,
  \item retained-teacher accelerations that predict a solver-relevant latent state but still end in a corrective teacher solve,
  \item positive-feature and low-rank coupling routes only as explicitly approximate or semantic-changing candidates,
  \item direct learned maps such as ICNN/Brenier/Monge-gap below those for immediate LEDH-PFPF-OT work,
  \item dynamic and operator methods as larger re-architecture ideas rather than first practical targets.
\end{enumerate}

\paragraph{Why exact-semantics streaming/online/GPU Sinkhorn is first.}
This route is ranked first because it changes the engineering realization while
preserving the mathematical object.  In the entropic OT chapter the teacher solve is
still the same finite entropic problem, with the same scaling equations and the same
barycentric transport output, even if the implementation is chunked, streamed, or
GPU-optimized.  So the main question is numerical validity under the existing contract,
not semantic reinterpretation.  This makes it the strongest reference lane for parity,
residual checks, and downstream gradient diagnostics.

\paragraph{Why retained-teacher accelerations are next.}
The retained-teacher chapter gives the cleanest neural/scalable extension of that
reference lane.  Meta OT and UNOT are strong here because they do not ask BayesFilter
to trust a learned final transport answer.  They predict a teacher-native latent state
(dual or solver coordinate), then run a corrective teacher solve.  Mathematically,
this preserves the distinction
\begin{equation}
\text{predicted latent state}
\neq
\text{final transport answer},
\label{eq:ledhpfpfot-retained-distinction}
\end{equation}
and therefore makes the first admissibility diagnostic very sharp: does warm start plus
correction preserve the teacher residual contract and barycentric output semantics?
That is why retained-teacher methods are the strongest practical next step after the
exact reference lane.

\paragraph{Why positive-feature and low-rank routes are lower.}
These routes are useful, but they should be ranked below retained-teacher methods
because they change the executed object more aggressively.  Positive-feature methods
can replace the Gibbs kernel by a feature approximation; low-rank coupling methods
restrict the feasible coupling family.  In both cases, the resulting transport layer
may still be useful, but it no longer inherits exact-semantics status automatically.
So these routes need their own evidence contracts.  Dense-reference mismatch is then an
explanatory diagnostic, not automatically a failure or a success.  That is why they are
later candidates rather than the default first recommendation.

\paragraph{Why direct learned maps are lower still.}
ICNN/Brenier and related direct-map methods ask LEDH-PFPF-OT to change the learned
object altogether.  Instead of preserving a retained entropic teacher, they try to
learn a final map or one of its surrogates.  That is a much larger step.  It changes
not only runtime and parameterization, but also what the implementation must verify at
deployment time.  In a finite-cloud BayesFilter setting, the Brenier logic is valuable
structural guidance, but it does not by itself justify replacing the current finite
entropic correction layer as the immediate next move.

\paragraph{Why dynamic and operator methods are last for immediate work.}
Dynamic and operator methods such as GeONet and its neighbors answer a broader
question: can we learn a whole path, a motion law, or a measure-to-measure operator?
That is scientifically important, but it is not the same question as improving the
current LEDH-PFPF-OT bottleneck.  These methods are therefore best viewed as future
re-architecture candidates rather than the first implementation target.

\paragraph{Concrete recommendation.}
If the research goal is to improve LEDH-PFPF-OT without reopening the whole transport
object, the first work items should be:
\begin{enumerate}[label=(\arabic*)]
  \item stabilize and optimize the exact entropic OT reference lane under the declared GPU/TF32/streaming target,
  \item implement retained-teacher warm-started acceleration, with Meta OT as the first code-backed source and UNOT as the strongest discrete-operator alternative,
  \item only after those pass teacher-preservation diagnostics, open explicit experiment plans for positive-feature and low-rank coupling replacements,
  \item defer direct learned maps and dynamic/operator re-architectures until the retained-teacher lane has either succeeded or clearly saturated.
\end{enumerate}

\paragraph{What this ranking does not claim.}
This ranking is about implementation priority for LEDH-PFPF-OT, not scientific
closure.  It does \emph{not} establish posterior correctness, HMC readiness, dense-
Sinkhorn equivalence for approximate routes, or production readiness from the current
diagnostics alone.  The precision result note already supports only a much weaker
claim: TF32 drift on the tested fixture is small relative to PF Monte Carlo noise,
but that is still not an HMC-level promotion argument.  So the section should be read
as a disciplined work-priority statement, not as a correctness certificate.

\section{What students should remember}

A first-year PhD student should leave this note with the following hierarchy in
mind.
\begin{enumerate}[label=(\arabic*)]
  \item OT-ICNN is an older Brenier-style baseline and was not disproven by the
  earlier BayesFilter benchmark.
  \item Wang teaches how conditional posterior transport can be written as a
  change-of-variables problem indexed by the observation.
  \item Al-Jarrah teaches how a filtering-native posterior-update map can be
  characterized by conditional expectation identities.
  \item Hosseini teaches the theorem-level conditional OT object that the more
  practical papers are approximating.
  \item Meta OT and UNOT teach how to keep a corrective teacher solve in the loop
  while learning to solve it faster.
\end{enumerate}

\section{Final conclusion}

The correct conclusion of the reread is therefore not a one-paper winner.  It is a
structured lesson about objects.

\begin{enumerate}[label=(\arabic*)]
  \item The old OT-ICNN lane should be closed as an unfaithful bridge benchmark,
  not as a verdict on the original Makkuva et al. method.
  \item The main modern scientific lane for BayesFilter is conditional posterior
  transport, led by Wang, Al-Jarrah, and Hosseini.
  \item The main conservative implementation lane is retained-teacher
  acceleration, led by Meta OT and UNOT.
  \item Dynamic and operator papers are important, but they answer a larger and
  harder question than BayesFilter immediately needs.
\end{enumerate}

That is the mathematically and pedagogically supported conclusion of the present
note.

\chapter*{Code snapshot summary}
\addcontentsline{toc}{chapter}{Code snapshot summary}

\subsection*{Meta OT}
\paragraph{GitHub / source.}
`https://github.com/facebookresearch/meta-ot`

\paragraph{Stack.}
JAX + OTT, with PyTorch also mentioned in the local README setup context.

\paragraph{Status / watch-outs.}
Official repo confirmed.  This is the strongest code-backed retained-teacher
starting point in our current notes, but students should still expect a real
adaptation burden when translating the teacher object and framework conventions
into BayesFilter.

\subsection*{UNOT}
\paragraph{GitHub / source.}
`https://github.com/GregorKornhardt/UNOT`

\paragraph{Stack.}
Python + PyTorch, with an FNO-style neural-operator setup and pretrained models.

\paragraph{Status / watch-outs.}
Official repo confirmed.  This is one of the strongest discrete entropic OT code
bases in the current survey, but the main watch-out is the representation bridge
from grid- or resolution-centric inputs to unordered particle clouds.

\subsection*{OT-ICNN / Makkuva}
\paragraph{GitHub / source.}
`https://github.com/AmirTag/OT-ICNN`

\paragraph{Stack.}
The local clone strongly suggests a Python / PyTorch-style codebase, though the
current repo notes do not yet summarize the framework as formally as they do for
Meta OT or UNOT.

\paragraph{Status / watch-outs.}
Official repo confirmed.  Best read as a faithful reproduction baseline for the
ICNN/Brenier route, not as a ready-made BayesFilter filtering module.  The main
watch-out is to preserve the original minimax / convex-potential contract rather
than collapse it into a generic direct regressor.

\subsection*{Wang conditional OT family}
\paragraph{GitHub / source.}
No clean confirmed local repository is recorded in the current ledger for this
exact paper title, though the surrounding conditional OT / Bayesian inference
family is clearly code-adjacent rather than theory-only.

\paragraph{Stack.}
Not yet recorded in a repo-verified way in the current local notes.

\paragraph{Status / watch-outs.}
Students should expect substantial fresh implementation.  The main watch-out is
that the static inverse-map route and the dynamic flow route are not the same
computational object and should not be blended into one implementation by habit.

\subsection*{Al-Jarrah filtering OT}
\paragraph{GitHub / source.}
No confirmed public repo in the current local ledger.

\paragraph{Stack.}
Not yet recorded in the current notes.

\paragraph{Status / watch-outs.}
This is a high-value paper but currently reads as a fresh reconstruction route for
BayesFilter.  The main watch-out is not code reuse but semantic correctness:
getting the conditional map, sample bookkeeping, and deployed posterior-update
operator right.

\subsection*{Hosseini conditional OT on function spaces}
\paragraph{GitHub / source.}
No confirmed public repo in the current local ledger.

\paragraph{Stack.}
Not yet recorded in the current notes.

\paragraph{Status / watch-outs.}
The paper is theorem-backed but effectively paper-first in our local
implementation record.  The main watch-out is that a student must still choose a
finite surrogate carefully before there is anything concrete to port.

\subsection*{GeONet}
\paragraph{GitHub / source.}
No confirmed public repo in the current local ledger.

\paragraph{Stack.}
Not yet recorded in the current notes.

\paragraph{Status / watch-outs.}
Strong conceptual dynamic/operator anchor, but still effectively a fresh
implementation project in our notes.  The main watch-out is not code reuse but
managing the richer path/PDE object without collapsing it into a simple endpoint
map.

\subsection*{Other Chapter 5 frontier papers}
\paragraph{GitHub / source.}
Mixed or often unclear in the current local research artifacts.

\paragraph{Stack.}
Mixed, unconfirmed, or not yet recorded in current notes.

\paragraph{Status / watch-outs.}
Students should distinguish code existence from implementation fit.  Several of
these routes are conceptually useful, but many still look adaptation-heavy even
when some code may exist elsewhere.

This summary is intentionally conservative.  It reuses the repo's earlier
source-availability ledger, implementation-fit note, and now the locally cloned
repositories where available.  It should be read as an implementation-feasibility
checklist rather than as a definitive software-maintenance, license, or
reproducibility audit.

\chapter*{References}
\addcontentsline{toc}{chapter}{References}
\begin{thebibliography}{99}

\bibitem{makkuva2020}
A. V. Makkuva, A. Taghvaei, S. Oh, and J. Lee.
\newblock Optimal transport mapping via input convex neural networks.
\newblock In \emph{Proceedings of the 37th International Conference on Machine Learning}, pages 6672--6681, 2020.

\bibitem{bunne2022}
Charlotte Bunne, Andreas Krause, and Marco Cuturi.
\newblock Supervised training of conditional Monge maps.
\newblock In \emph{Advances in Neural Information Processing Systems}, 2022.

\bibitem{amos2023}
Brandon Amos, Giulia Luise, Samuel Cohen, and Ievgen Redko.
\newblock Meta optimal transport.
\newblock In \emph{Proceedings of the 40th International Conference on Machine Learning}, pages 1472--1509, 2023.

\bibitem{scetbon2023}
Meyer Scetbon, Michal Klein, Giovanni Palla, and Marco Cuturi.
\newblock Unbalanced low-rank optimal transport solvers.
\newblock arXiv:2305.19727, 2023.

\bibitem{kassraie2024}
Parnian Kassraie, Aram-Alexandre Pooladian, Michal Klein, James Thornton, Jonathan Niles-Weed, and Marco Cuturi.
\newblock Progressive entropic optimal transport solvers.
\newblock arXiv:2406.05061, 2024.

\bibitem{gazdieva2024}
Milena Gazdieva, Arip Asadulaev, Evgeny Burnaev, and Alexander Korotin.
\newblock Light unbalanced optimal transport.
\newblock In \emph{Advances in Neural Information Processing Systems}, 2024.

\bibitem{geonet2024}
Andrew Gracyk and Xiaohui Chen.
\newblock GeONet: a neural operator for learning the Wasserstein geodesic.
\newblock In \emph{Proceedings of the Fortieth Conference on Uncertainty in Artificial Intelligence}, pages 1638--1648, 2024.

\bibitem{wang2025}
Zheyu Oliver Wang, Ricardo Baptista, Youssef Marzouk, Lars Ruthotto, and Deepanshu Verma.
\newblock Efficient neural network approaches for conditional optimal transport with applications in Bayesian inference.
\newblock arXiv:2310.16975, revised 2025.

\bibitem{aljarrah2025}
Mohammad Al-Jarrah, Niyizhen Jin, Bamdad Hosseini, and Amirhossein Taghvaei.
\newblock Nonlinear filtering with Brenier optimal transport maps.
\newblock In \emph{Proceedings of the 41st International Conference on Machine Learning}, 2024.

\bibitem{hosseini2025}
Bamdad Hosseini, Alexander W. Hsu, and Amirhossein Taghvaei.
\newblock Conditional optimal transport on function spaces.
\newblock arXiv:2311.05672, 2023.

\bibitem{geuter2025}
Jonathan Geuter, Gregor Kornhardt, Ingimar Tomasson, and Vaios Laschos.
\newblock Universal neural optimal transport.
\newblock In \emph{Proceedings of the 42nd International Conference on Machine Learning}, 2025.

\bibitem{chen2025}
Peter Chen, Yue Xie, and Qingpeng Zhang.
\newblock Displacement-sparse neural optimal transport.
\newblock arXiv:2502.01889, 2025.

\bibitem{chaudhari2025}
Shreyas Chaudhari, Srinivasa Pranav, and Jos\'e M. F. Moura.
\newblock GradNetOT: learning optimal transport maps with GradNets.
\newblock arXiv:2507.13191, 2025.

\bibitem{chenliu2025}
Chen and Liu.
\newblock Dynamic optimal transport and learned motion fields.
\newblock Working title placeholder; exact citation metadata still to be confirmed.

\bibitem{girshfeld2025}
Inga Girshfeld and Xiaohui Chen.
\newblock Neural local Wasserstein regression.
\newblock In \emph{Proceedings of the 1st Conference on Topology, Algebra, and Geometry in Data Science}, pages 79--89, 2026.

\bibitem{dumont2026}
Th\'eo Dumont, Th\'eo Lacombe, and Fran\c{c}ois-Xavier Vialard.
\newblock Learning Monge maps by lifting and constraining Wasserstein gradient flows.
\newblock arXiv:2603.25182, 2026.

\bibitem{park2026}
Yesom Park, Eric Gelphman, Stanley Osher, and Samy Wu Fung.
\newblock Fixed-point neural optimal transport without implicit differentiation.
\newblock arXiv:2605.10792, 2026.

\bibitem{micheli2026}
Alessandro Micheli, Yueqi Cao, Anthea Monod, and Samir Bhatt.
\newblock Riemannian neural optimal transport.
\newblock arXiv:2602.03566, 2026.

\bibitem{vandergrift2026}
Matthew Vandergrift, Martha White, Yury Polyanskiy, Philippe Rigollet, and Lazar Atanackovic.
\newblock Measure-to-measure regression with transformers.
\newblock arXiv:2605.28075, 2026.

\end{thebibliography}

\end{document}
